In Section~\ref{sec:infproof}, a proof of Theorem~\ref{th:upper} is
presented; the proofs of some auxiliary results are deferred to
Sections~\ref{subsec:dead} and~\ref{subsec:virtrevers}. 
We begin with several
observations and conventions.


\paragraph{1-exp and 2-exp functions.}

We have defined 1-exp and 2-exp functions as functions $f\colon\setN^2\to\setN$
satisfying
\begin{equation*}
	f(m,i)\leq (2+m)^{\text{poly}(i)}\text{ and }
f(m,i)\leq (2+m)^{2^{\text{poly}(i)}},
\end{equation*}
	respectively.
For technical convenience, we implicitly assume that each such 
function $f$ that we consider is \emph{non-decreasing} in both arguments; that is, 
$m \leq m'$ and $i \leq i'$ imply $f(m, i) \leq f(m', i')$. 
Furthermore, we assume that $f$ is \emph{extensive}, meaning that $m \leq f(m,i)$.

It is straightforward to verify that the class of 1-exp functions, as
well as  the class of 2-exp functions, 
 is
closed under standard operations such as sum, product, and
composition, where the latter is
defined as \[(f\circ g)(m,i)=f(g(m,i),i).\]
For instance, if $f(m,i)\leq	(2+m)^{{p_1(i)}}$
and 
 $g(m,i)\leq	(2+m)^{{p_2(i)}}$ with $p_1(i)\geq 1$ and
 $p_2(i)\geq 1$, then  
\[(f\circ g)(m,i)\leq
\left(2+(2+m)^{{p_2(i)}}\right)^{{p_1(i)}}\leq 
(2+m)^{{2\cdot p_2(i)\cdot p_1(i)}}.\]
The functions $p_1, p_2$
can be, in particular, polynomials or exponential functions.

 We define the $j$-th iterate of $f$ recursively
as follows:
\begin{equation*}
	f^{(0)}(m,i)=m,\text{ and }
f^{(j+1)}(m,i) = f(f^{(j)}(m,i), i).
\end{equation*}
Note that 
if $f(m,i)$ is a 1-exp function or a 2-exp function, then $f'(m,i)=f^{(i)}(m,i)$ 
is a 2-exp function. Indeed,
if $f(m,i)\leq (2+m)^{{p(i)}}$, where $p$ is a polynomial or
exponential function, then
\[f'(m,i)=f^{(i)}(m,i)\leq (2+m)^{2^{i-1}\cdot
(p(i))^i}=(2+m)^{2^{(i-1)+ i\cdot \log
p(i)}}\leq(2+m)^{2^{\text{poly}(i)}}.\]
We use these closure properties
implicitly when deriving the existence of specific 1-exp or 2-exp functions 
throughout this section.

\paragraph{1-exp and 2-exp constants.}

For convenience, when considering a $d$-dim net $A$, we also
refer to \emph{1-exp} and \emph{2-exp constants}. We say that $C$ is a 1-exp (resp.~2-exp)
constant related to $A$ if
$C=f(\nnorm{A},d)$ for some 1-exp (resp.~2-exp) function that is independent of $A$.
For instance, we will use a 2-exp constant $C\dead^A$ associated
with a 2-exp function $f\dead$; the superscript $A$ clarifies that 
$C\dead^A=f\dead(\nnorm{A},d)$, where $d$ is the dimension of the net $A$.

\subsection{Proof of
Theorem~\ref{th:upper}.}\label{sec:infproof}
\hfill\\

We aim to show the existence of a 2-exp function $f$ such that for every $d$-dim
net $A$ possessing a live bottom SCC $X\subseteq\setN^d$ (where every
$x\in X$ is a live marking and $\Reach(x)=X$ for all
$x\in X$),
 there exists a live marking
$x\in\setN^d$ satisfying $\nnorm{x}\leq f(\nnorm{A},d)$.
 Recall that 
such a net $A$ is reversible by Lemma~\ref{lem:consstrbound}.

In this section, we provide the overall proof, deferring certain
technical parts 
to the subsequent sections.

\begin{rem}\label{rem:Lerresults}
Some technical issues are addressed by referring to results from~\cite{DBLP:conf/fsttcs/Leroux19}. Specifically, these concern
$d$-dimensional \emph{extractors} (used to distinguish between small
	and large components of markings, or of sets of markings),
a 1-exp \emph{virtual reachability} function $f\vr$ for reversible
	Petri nets with states (PNSs), and 
	a 2-exp \emph{mutual-reachability} function $f\mr$ for general
	Petri nets.
\end{rem}

\paragraph{Dead markings, 2-exp constant $\bm{C\bdead^A}$}
Recall that a marking $x\in\setN^d$ of a $d$-dim net $A$ is live if and only if there is no
dead marking in $\Reach(x)$
(Proposition~\ref{prop:nonlivereachdead}). By our convention, a
marking $y$ is dead if at least one action $a$ is dead at $y$, that
is, $y'\not\geq a_-$ for all $y'\in\Reach(y)$.
Our goal is to bound the minimal size of solutions $x\in\setN^d$ to the constraint
\begin{equation}\label{eq:firstconstr}
	\textnormal{``for each dead marking $y$, we have  $x\not\reach
	y$''.}
\end{equation}
In Section~\ref{subsec:dead} we recall, by
Lemma~\ref{lem:rackoffdead} using  Rackoff's
result in~\cite{DBLP:journals/tcs/Rackoff78}, that 
\begin{equation*}
	\textnormal{\emph{the set of dead markings} of a $d$-dim net
	$A$} 
\end{equation*}
	can be presented
as the set of solutions $y\in\setN^d$ 
of a~disjunction
\begin{equation}\label{eq:alldead}
	\bigvee_{I\subseteq [1,d]} \,\bigvee_{j=1}^{k_I}
	\left(\bvar{y}\restr{I}=\alpha_{(I,j)}\right)
\end{equation}
where 
each (restricted) marking $\alpha_{(I,j)}\in\setN^{I}$ is \emph{dead in the 
restricted net} $A\restr{I}$ and satisfies 
\[\nnorm{\alpha_{(I,j)}}< C\dead^A,\] for some 2-exp
constant $C\dead^A$.
We allow $k_I=0$ for some $I\subseteq [1,d]$; the corresponding
disjunction $\bigvee_{j=1}^{0}$ is empty, and thus false.  The
disjunction~\eqref{eq:alldead} can therefore be restricted to
$\bigvee_{I\in\mathcal{D}}\cdots$, where
$\mathcal{D}=\{I\subseteq[1,d] \mid k_I>0\}$.  
Note that $\emptyset\notin\mathcal{D}$: in the restricted net
$A\restr{\emptyset}$ no action is ever dead (recall
Remark~\ref{rem:dimzero}), hence no dead $\alpha_{(\emptyset,j)}$
exists.


\begin{rem}
By our convention, the expression ``some 2-exp constant $C\dead^A$''
	means that there exists a non-decreasing function $f\dead(m,i)\leq
	(2+m)^{2^{\text{poly}(i)}}$
such that $C\dead^A=f\dead(\nnorm{A},d)$ for
any $d$-dim net $A$.
A more detailed analysis would involve stating
that	$\nnorm{\alpha_{(I,j)}}< f\dead(\nnorm{A},|I|)$, where 
	$f\dead(\nnorm{A},1)=\nnorm{A}$, and discussing the tight
	relation to the coverability problem (for markings $a_-$). 
	However, for our  purposes, it suffices to simplify the
	analysis by referring only to
	$C\dead^A$, even when a restricted net $A\restr{I}$ is
	considered.

	Technically, we do not rely on the fact that the set of dead
	markings of $A$ is downward closed (that is, if $y$ is dead and $y'\leq y$,
	then $y'$ is also dead). We only use the property that,
	given a marking $y$ of a net $A$, 
if $y\restr{I}$ is dead in the
	restricted net $A\restr{I}$, then $y$ is dead in $A$
	(Proposition~\ref{prop:restrdead}).
\end{rem}

\paragraph{(Virtual) reachability, 1-exp
$\bm{f\bvr}$ and $\bm{C\bvr^A}$, 2-exp $\bm{C\bsuf^A=\max\{C\bdead^A,C\bvr^A\}}$.}
Our proof of Theorem~\ref{th:upper} is based on replacing
the non-reachability relation $\not\reach$ in 
constraint~(\ref{eq:firstconstr}) with the virtual reachability relation
$\vreach$, in a way that will be clarified later. 
Markings in which all components are sufficiently large will play a
crucial role.

\begin{defi}[\textbf{(Sub)markings with all components above a level,
	$\bm{x\in\upC{C}}$}]
\hfill\\
We say that a \emph{vector} $x\in\setN^d$ is \emph{above} a
\emph{level} $C\in\setN$, or simply
$x\in\setN^d$ is \emph{above} $C\in\setN$, denoted
 $x\in\upC{C}$,	if $x(i)\geq C$ for all $i\in[1,d]$.
	A restricted vector
	$x\restr{J}$ with 
	$J\subseteq[1,d]$ is above $C$,
	denoted $x\restr{J}\in\upC{C}$, if  $x(i)\geq C$ for all $i\in
	J$.
\end{defi}

% Fix Indentation
\noindent
The following lemma is a consequence of~\cite[Lemma~5]{DBLP:conf/fsttcs/Leroux19}.

\begin{lem}[\textbf{Virtual and standard reachability in
	reversible nets, $\bm{f\bvr}$,
	$\bm{C\bvr^A}$}]\label{lem:fvr}
\hfill\\	
	There exists a 1-exp function $f\vr$ such that 
for every reversible $d$-dim net $A$ we have:
	\begin{enumerate}
		\item
If $x\vreach y$, then  $x\virtual{\sigma}y$
for some $\sigma$ satisfying
			$|\sigma|\leq
			f\vr(\nnorm{A},d)\cdot\nnorm{y-x}=
		C\vr^A\cdot \nnorm{y-x}$.
\item
If $x\vreach y$
			and  $x,y\in\upC{C\vr^A}$, 
	then $x\step{\sigma}y$
	for some $\sigma$ satisfying
$|\sigma|\leq
		C\vr^A\cdot \nnorm{y-x}$.
	\end{enumerate}
	\end{lem}
% Fix Indentation
\noindent
Intuitively, this can be seen by recalling that $x\vreach y$ means
$y-x\in A\D^*$. 
We highlight the fact that, given a reversible $d$-dim net $A$,  virtual reachability coincides with
standard reachability above $C\vr^A$:
\begin{equation*}
	x\vreach y\iff x\reach y,\text{ for all
	}x,y\in\setN^d\cap \upC{C\vr^A}.
\end{equation*}
For later convenience, we define a sufficiently large 2-exp
constant:
\begin{equation}\label{eq:Csuf}
C\suf^A=\max\{C\dead^A,C\vr^A\}.
\end{equation}


\paragraph{Segment decomposition of virtual executions in reversible nets,
1-exp $\bm{C\bshift^A}$.}
\ \\
In Section~\ref{subsec:virtrevers}, 
we show that
if $x\virtual{*}y$ in a \emph{reversible} $d$-dim net $A$, then
there exists
a~virtual execution 
\begin{equation}\label{eq:virtsegment}
x=x_0\virtual{*}x_1\virtual{*}x_2\cdots\virtual{*}x_k=y
\end{equation}
consisting of \emph{1-exp segments} 
that \emph{stepwise approach the
target} $y$ from $x$.
The segments $x_j\vreach x_{j+1}$ are 1-exp bounded in the sense that
\begin{equation}\label{eq:Cshift}
0<\nnorm{x_{j+1}-x_j}\leq C\shift^A
\end{equation}
for some 1-exp constant $C\shift^A$ (satisfying
$C\shift^A=f\shift(\nnorm{A},d)\leq
(2+\nnorm{A})^{\text{poly}(d)}$).
By the term \emph{stepwise approaching},
	we mean that the sequence $x = x_0, x_1, \dots, x_k = y$ satisfies:
\begin{itemize}
    \item For each $i \in [1, d]$ with $x(i) \leq y(i)$, there exists
	    $j_i\in [0,k]$ such that
		\[ x(i) < x_1(i) < \cdots < x_{j_i}(i) = x_{{j_i}+1}(i) = \cdots = x_k(i)
		= y(i); \]
    \item For each $i \in [1, d]$ with $x(i) > y(i)$,  there exists
	    $j_i\in [0,k]$ such that
		\[ x(i) > x_1(i) > \cdots > x_{j_i}(i) = x_{{j_i}+1}(i) = \cdots = x_k(i)
		= y(i). \]
\end{itemize}

% Fix Indentation
\noindent
Note that Lemma~\ref{lem:fvr}(1) implies that for each segment $x_j\vreach x_{j+1}$
of~(\ref{eq:virtsegment}) there exists  $\sigma$ with
$|\sigma|\leq C\vr^A\cdot C\shift^A$ such that
$x_j\virtual{\sigma}x_{j+1}$,
and even $x_j\step{\sigma}x_{j+1}$ if $x_j,x_{j+1}\in\upC{C\vr^A}$.


\paragraph{Small and large components of vectors, relative to
extractors.}
We need to distinguish between relatively small and large components 
of vectors $x \in \setN^d$ (markings of a net $A$). 
To identify small components, we employ the notion of 
extractors, following~\cite{DBLP:conf/fsttcs/Leroux19}.


\begin{defi}[\textbf{Extractors $\bm{\lambda}$	of dimension $\bm{d}$}]\label{def:extract}
\hfill\\	
For $d\geq 1$, by a \emph{$d$-dim extractor} $\lambda$
we mean a tuple 
	\[(\lambda_1,\lambda_2,\dots,\lambda_d)\in(\setN_+)^d
	\text{ where }\lambda_1\leq\lambda_2 \leq \dots \leq
	\lambda_{d}.\]
For technical convenience, we also refer to the values $\lambda_0$
	and $\lambda_{d+1}$, by
implicitly setting $\lambda_0=1$ and $\lambda_{d+1}=\lambda_d$.
\end{defi}

% Fix Indentation
\noindent
The following claim is straightforward (for further details
see~\cite{DBLP:conf/fsttcs/Leroux19}).

\begin{prop}[\textbf{The set $\bm{I^\lambda_x\subseteq[1,d]}$ of
	(indices of) small
	components of $\bm{x\in\setN^d}$}]\label{prop:Ilambdax}
\hfill\\	
Given a $d$-dim extractor $\lambda$,
for each $x\in\setN^{d}$ there exists a unique maximal
	set $I_x^\lambda\subseteq [1,d]$ such that 
	\[\bigwedge_{i\in
	I^\lambda_x}\left(x(i)<\lambda_{|I^\lambda_x|}\right)\ \land\ 
	\bigwedge_{i\in[1,d]\smallsetminus
	I^\lambda_x}\left(x(i)\geq\lambda_{|I^\lambda_x|+1}\right).
	\]
	More generally, for each subset $J\subseteq [1,d]$, there is 	a unique maximal
	set $I_x^{(\lambda,J)}\subseteq J$ such that 
	\[\bigwedge_{i\in
	I^{(\lambda,J)}_x}\left(x(i)<\lambda_{|I^{(\lambda,J)}_x|}\right)\
	\land \
	\bigwedge_{i\in J\smallsetminus
	I^{(\lambda,J)}_x}\left(x(i)\geq\lambda_{|I^{(\lambda,J)}_x|+1}\right).
	\]
\end{prop}	
% Fix Indentation
\noindent
For instance, if $\lambda=(3,17,258,1500)$ and $x=(11,1600,2,300)$,
$x'=(16,5000,15,5000)$, then $I^\lambda_x=I^\lambda_{x'}=\{1,3\}$.
The values $x(i)$ of the  components $i\in I^\lambda_x$ are viewed as
\emph{small}, while the values $x(i)$ for $i\in [1,d]\smallsetminus I^\lambda_x$ are
\emph{large}. The gap between the values of large components and of
small components of $x$ is bigger
than $\lambda_{j+1}-\lambda_j$, where $j=|I^\lambda_x|$.
In our example, with $|I^\lambda_x|=2$, the gap is bigger than $\lambda_3-\lambda_2=258-17=241$.
\begin{figure}[t]	
\begin{tikzpicture}
    \begin{scope}
        \draw[black, thick] (0,0) -- (0,7);
        \draw[black, thick] (0,0) -- (3,0);
        \draw[black, thick] (3,0) -- (3,7);
        \node[anchor=south,rotate=90] at (0,3.5) {Number of tokens};

        \draw[dashed, black] (0.98,0) -- (0.98,7);
        \node[anchor=north] at (0.5,0) {$I$};
        \node[anchor=north] at (2,0) {$[1,d] \smallsetminus I$};
        
        \draw[densely dotted, black] (0,2) -- (3,2);
        \node[anchor=west] at (3,2) {$\lambda_{\sizeof{I}}$};
        
        \draw[densely dotted, black] (0,5.5) -- (3,5.5);
        \node[anchor=west] at (3,5.5) {$\lambda_{\sizeof{I}+1}$};
        \draw[densely dotted, black] (0,0.75) -- (3,0.75);
        \node[anchor=west] at (3,0.75) {$C\suf^A$};

        \drawconf{blue}{very thick}{
            1.5,1.4,1.1,0.9,1.9,
            7.0,6.5,5.5,5.8,6.1,
            5.9,5.7,5.8,5.9,6.4
        }{0.2}{0}
	    \node[anchor=west,text=blue] at (3,6.4) {$\bm{x}$};
    \end{scope}
    \begin{scope}[xshift=7cm,yshift=0cm]
        \draw[black, thick] (0,0) -- (0,7);
        \draw[black, thick] (0,0) -- (3,0);
        \draw[black, thick] (3,0) -- (3,7);
        \node[anchor=south,rotate=90] at (0,3.5) { Number of tokens};

        \draw[dashed, black] (0.98,0) -- (0.98,7);
        \node[anchor=north] at (0.5,0) {$I$};
        \node[anchor=north] at (2,0) {$[1,d] \smallsetminus I$};
        
        \draw[densely dotted, black] (0,2) -- (3.2,2);
        \node[anchor=west] at (3.2,2) {$\lambda_{\sizeof{I}}$};
        \draw[black, very thick] (3.2,1.9) -- (3.2,2.8);
        \draw[densely dotted, black] (0,2.8) -- (3.2,2.8);
	    \node[anchor=west] at (3.2,2.4)
	    {$x'(i)-(\lambda_{\sizeof{I}}-1)>0$};
	    \node[anchor=west] at (3.7,3.0)
	    {$i\in I$};
        
        \draw[densely dotted, black] (0,4.5) -- (3.2,4.5);
        \node[anchor=west] at (3.2,5.0)
	    {$\lambda_{\sizeof{I}+1}-x'(i)>0$};
	 \node[anchor=west] at (5.1,4.5)
	    {$i\notin I$};
        \draw[black, very thick] (3.2,4.5) -- (3.2,5.5);
        \draw[densely dotted, black] (0,5.5) -- (3.2,5.5);
        \node[anchor=west] at (3.2,5.5) {$\lambda_{\sizeof{I}+1}$};
        \draw[densely dotted, black] (0,0.75) -- (3,0.75);
        \node[anchor=west] at (3,0.75) {$C\suf^A$};

         \drawconf{red}{very thick}{
	    1.5,1.0,0.8,2.8,2.3,
            7.0,6.5,6.1,6.2,4.9,
            4.5,4.8,5.1,6.5,6.3
        }{0.2}{0}
	    \node[anchor=west,text=red] at (3,6.3) {$\bm{x'}$};
    \end{scope}
\end{tikzpicture}
	
	\caption{Index set $I=I^\lambda_x$ (left);
	depiction of $\dist{\lambda}{I}{x'}$ (right).}\label{fig:extractor}
\end{figure}

Figure~\ref{fig:extractor} (left) depicts a marking $x\in\setN^d$
and highlights the index set $I=I^\lambda_x\subseteq[1,d]$. 
The marking $x$ is also an example of a marking 
above the level $C\suf^A$ (defined in~\eqref{eq:Csuf});
such markings are particularly important in our later reasoning.

Note that $I^\lambda_x=[1,d]$ if and only if $x(i)< \lambda_d$ for
every $i\in[1,d]$;
the condition  $x(i)\geq \lambda_{d+1}$ holds vacuously
for all $i\in\emptyset$.
If $J'=\{i\in[1,d]\mid x(i)\geq \lambda_d\}$, then 
$I^\lambda_x=I^{(\lambda,J)}_x$ for $J=[1,d]\smallsetminus J'$.
Specifically, if $J'=[1,d]$, then $I^\lambda_x=\emptyset$; the
condition
$x(i)<\lambda_{0}$ holds vacuously for all $i\in\emptyset$.
In general,  $I^\lambda_x=\emptyset$ if and only if for each
$\emptyset\neq I\subseteq[1,d]$, there exists $i\in I$ such that
$x(i)\geq \lambda_{|I|}$.

Later, we apply \emph{2-exp $d$-dim extractors} $\lambda$ to the $d$-dim nets $A$ under
consideration, meaning 
that 
\begin{equation*}
\lambda_d=f(\nnorm{A},d)\leq(2+\nnorm{A})^{2^{\text{poly}(d)}}
\end{equation*}
for some (implicit) 2-exp function $f(m,i)$ which is independent of $A$.
We use such extractors for which
the gaps $\lambda_{j+1}-\lambda_j$ between the large and small
components grow sufficiently with the index $j$.
However, before defining specific extractors, we introduce certain linear systems with
2-exp solutions.



\paragraph{Linear systems underlying the proof of Theorem~\ref{th:upper}.}
In Definition~\ref{def:crucialLS} we introduce a class of linear systems
on which our proof rests; Proposition~\ref{prop:Affsmall} shows that
these systems have 2-exp solutions.

The definition uses the notion of the \emph{distance of a vector to a
$\lambda$-case}, denoted $\dist{\lambda}{I}{x}$ for a vector
$x\in\setN^d$ and a $\lambda$-case $I$, where
$\emptyset\neq I\subseteq[1,d]$.
This distance is a nonnegative integer, and it is zero whenever
$I^\lambda_x=I$.
See Figure~\ref{fig:extractor} (right) for an illustration.


\begin{defi}[\textbf{Distance $\bm{\dist{\lambda}{I}{x}}$ of $\bm{x\in\setN^d}$ to the
	$\bm{\lambda}$-case $\bm{\emptyset\neq I\subseteq[1,d]}$}]
\hfill\\
Given a $d$-dim extractor $\lambda$, the \emph{distance of}
	$x\in\setN^d$ \emph{to} a \emph{$\lambda$-case} $I$, where
	$\emptyset\neq I\subseteq
[1,d]$, is 
	\[\dist{\lambda}{I}{x}=\max\left(\{0\}\cup
	\{x(i)-(\lambda_{|I|}-1)\mid i\in
I\}\cup\{\lambda_{|I|+1}-x(i)\mid i\in[1,d]\smallsetminus I\}\right).\]
\end{defi}	

% Fix Indentation
\noindent
The value $\dist{\lambda}{I}{x}$ is positive if
and only if there exists
$i\in [1,d]$ such that $i\in I$ and $x(i)\geq \lambda_{|I|}$, or
$i\in[1,d]\smallsetminus I$ and $x(i)<\lambda_{|I|+1}$.
For instance, if $\lambda=(3,17,258,1500)$ and $x=(11,1600,2,300)$,
then $\dist{\lambda}{\{1,3\}}{x}=0$, 
and
$\dist{\lambda}{\{1,2\}}{x}=\max(\{0\}\cup\{-5,1584\}\cup\{256,-42\})=1584$. 

\smallskip

The following definition of $(A,x_0,\lambda)$-linear systems aims at
transforming
 constraint~(\ref{eq:firstconstr}). 
For  \emph{reversible} $d$-dim nets $A$ and live 
markings $x_0$, 
our later choice of linear systems $S$
in this class will guarantee that 
every solution of $S$ presents
a nonempty set of at most $2^d$ \emph{live} markings of $A$ that are mutually \emph{virtually} reachable; in
fact, they are also mutually reachable since they all are above $C\vr^A$.
Hence, the non-reachability of dead markings in 
constraint~(\ref{eq:firstconstr}) is replaced with the positive
condition of virtual reachability of a specific subset of markings. 
The definition uses the 1-exp constant $C\shift^A$~(\ref{eq:Cshift}) and the 2-exp
constant $C\suf^A$~(\ref{eq:Csuf}) 
related to a $d$-dim
net $A$.

\begin{defi}[\textbf{$\bm{(A,x_0,\lambda)}$-linear systems for nets,
	markings, and 2-exp extractors}]\label{def:crucialLS}
Given a reversible $d$-dim net $A$ and a 2-exp extractor
	$\lambda=(\lambda_1,\lambda_2,\dots,\lambda_d)$,
	we define
the 2-exp bound	\begin{equation}\label{eq:BAffnets}
	B=B_{(A,\lambda)}=\max\{C\suf^A,\lambda_d+C\shift^A\}.
	\end{equation}
	Given a marking $x_0\in\setN^d\cap\upC{C\suf^A}$ (that is, 
	$\bigwedge_{i=1}^d x_0(i)\geq C\suf^A$), for each $\emptyset\neq I\subseteq[1,d]$ let 
	\[x^\close_I\in \setN^d\cap\upC{C\suf^A}\]
	be a fixed marking in $\Reach(x_0)$
	that is above $C\suf^A$ and minimizes the distance to
	the $\lambda$-case $I$, that is, $\dist{\lambda}{I}{x^\close_I}$ is
	minimal. 

Then the following 
linear system $S$ of dimension $d\cdot 2^d$, with
	a $d$-dim variable
	$\bvar{x}=(\bvar{x}_1,\bvar{x}_2,\dots,\bvar{x}_d)$ and $2^d-1$ $d$-dim variables 
	$\bvar{x}_I$, one for each $\emptyset\neq I\subseteq[1,d]$,
	is 
	\begin{center}
		an \emph{$(A,x_0,\lambda)$-linear system}:
	\end{center}		
\begin{equation}\label{eq:AxLS}
	(\bvar{x}\equiv_B x_0)
	\land
	\bigwedge_{\emptyset\neq I\subseteq[1,d]}\big((\bvar{x}_I\equiv_B x^\close_I)\land
	(\bvar{x}\virtual{*}\bvar{x}_I)\big),
\end{equation}
	where $\bvar{y}\equiv_B x$ is a shorthand for 

	\[\left(\bigwedge_{i\in[1,d], x(i)\leq B}
	(\bvar{y}_i=x(i))\right)\land 
	\left(\bigwedge_{i\in[1,d], x(i)> B} (\bvar{y}_i\geq
	B+1)\right).\]
\end{defi}	

\begin{prop}[\textbf{$\bm{(A,x_0,\lambda)}$-linear systems have 
	2-exp	solutions}]\label{prop:Affsmall}
\hfill\\
For every $(A,x_0,\lambda)$-linear system where $A$ has 
 dimension $d$,
there exists a solution 
	\begin{center}
	$x\in\setN^{d\cdot 2^d}$ such that 
	$\nnorm{x}\leq (2+\nnorm{A})^{2^{\text{poly}(d)}}$.
	\end{center}
	\end{prop}	
\begin{proof} 
Let $S$ be an $(A,x_0,\lambda)$-linear system where  $A$ has dimension $d$.
The system $S$ is
satisfiable, since the tuple $(x_0,(x^\close_I)_{\emptyset\neq I\subseteq [1,d]})$ is one
of its solutions.
	Due to Theorem~\ref{thm:smallsolutions},
there exists a solution
	$x\in\sem{S}$ such that
	\[\nnorm{x}\leq
		\mlcm{S} \cdot \left(1+ (\bar{d}+1)\,!\cdot (\bar{d}\cdot
		\nnorm{S}+1)^{\bar{d}}\right)\]
	where 	$\bar{d}=d\cdot 2^d$.
Recalling~(\ref{eq:AxLS})
and Corollary~\ref{cor:virtual2system}, we deduce that
	$\mlcm{S}\leq d\,!\cdot\nnorm{A}^d$, and $\nnorm{S}\leq
	1+B_{(A,\lambda)}\leq 
	(2+\nnorm{A})^{2^{\text{poly}(d)}}$. 
We thus obtain
\[\nnorm{x}\leq d\,!\cdot\nnorm{A}^d\cdot
 \left(1+ (d\cdot 2^d+1)\,!\cdot (d\cdot 2^d\cdot
		\nnorm{S}+1)^{d\cdot 2^d}\right),\]
which implies that 
	$\nnorm{x}\leq(2+\nnorm{A})^{2^{\text{poly}(d)}}$.
\end{proof}

% Fix Indentation
\noindent
Our goal now is to show certain \emph{lower bounds} for 
the increasing gaps
$\lambda_{j+1}-\lambda_j$ that ensure
that \emph{every solution} 
of an $(A,x_0,\lambda)$-linear system $S$
constitutes \emph{a set
of live markings} of $A$, under the condition that $x_0$ is a sufficiently
large \emph{live} marking.
Under this condition, the solution $(x_0,(x^\close_I)_{\emptyset\neq I\subseteq[1,d]})$ 
has the required property, since all markings $x^\close_I$ are reachable from
the live marking $x_0$, and are thus themselves live (by
Proposition~\ref{prop:liveunderreach}).
Once we establish that \emph{all solutions}
$(\bar{x}_0,(\bar{x}_I)_{\emptyset\neq I\subseteq[1,d]})$ of $S$ 
possess this property, we use the fact 
that the minimal solutions are guaranteed to be at most 2-exp by
Proposition~\ref{prop:Affsmall}.

This goal is achieved via Lemma~\ref{lem:finalbasicnets}, which builds on the quasi-dead marking property of \emph{reversible} nets established in Lemma~\ref{lem:reachquasi}.

\paragraph{Quasi-dead markings.}
We now show that, given a \emph{reversible} $d$-dim net $A$,
every nonlive marking $x_0\in\upC{\max\{C\suf^A,\nnorm{A}\cdot C\qd^A\}}$,
where $C\qd^A$ is a certain 2-exp constant,  can reach a dead marking via 
a  \emph{quasi-dead} marking $x\in\upC{C\suf^A}$.
By definition, quasi-dead markings must be above the level $C\suf^A$
and reach dead markings by executions
of length at most $C\qd^A$.
We use Figure~\ref{fig:markingstwo} for illustration.

\begin{defi}[\textbf{Quasi-dead markings, 2-exp function $\bm{f\bqd}$ and
constant	$\bm{C\bqd^A}$}]
\hfill\\	
For a 2-exp function $f\qd$ and a $d$-dim net $A$, a marking $x\in\setN^d$
	is \emph{$f\qd$-quasi-dead}, or simply \emph{quasi-dead} when $f\qd$
	is clear from context, if $x\in\upC{C\suf^A}$ 
and there exists an execution $x\step{\sigma}y$ such that $y$ is a dead
marking and $|\sigma|\leq
f\qd(\nnorm{A},d)=C\qd^A$.
\end{defi}

Lemma~\ref{lem:reachquasi} uses
Proposition~\ref{prop:simpletrick},
which is based on the following involved result~\cite[Theorem
2]{DBLP:conf/fsttcs/Leroux19}; 
 the result applies to general nets,
without assuming reversibility:
\begin{lemC}[\cite{DBLP:conf/fsttcs/Leroux19} {(\textbf{2-exp mutual
	reachability function $\bm{f\bmr}$, 2-exp constant
	$\bm{C\bmr^A}$})}]\label{lem:fmr}
	There exists a 2-exp function $f\mr$
	with the following property:
	For every $d$-dim net $A$, if $x\reach y$ and $y\reach x$, then
there exist action
	sequences $\sigma_1,\sigma_2$ such that 
	$x\step{\sigma_1}y\step{\sigma_2}x$ and
	\begin{equation*}
		|\sigma_1\sigma_2|\leq
		f\mr(\nnorm{A},d)\cdot\nnorm{x-y}=
		C\mr^A\cdot\nnorm{x-y}.
	\end{equation*}		
\end{lemC}

	\begin{prop}[\textbf{Down reachability $\bm{x\reach y}$ with
		$\bm{x\geq y}$}]\label{prop:simpletrick}
	\hfill\\
	Given a $d$-dim net $A$ and markings $x\geq y$ such that $x\reach y$, 
there exists a marking $x'\geq x$ satisfying the following conditions:
\begin{enumerate}[a)]	
	\item $\nnorm{x'-x}\leq C\mr^A\cdot \nnorm{x-y}$; 
		\item
			$x'\step{\sigma}y$ for some $\sigma$ 
			such that $|\sigma|\leq
			C\mr^A\cdot\nnorm{x-y}$;
		\item $x'(i)=x(i)$ if $x(i)=y(i)$
			(and $x'(i)\geq x(i)$ if $x(i)>y(i)$), for all
			$i\in[1,d]$.
\end{enumerate}
\end{prop}
\begin{proof}
Given a $d$-dim net $A$, and $x,y\in\setN^d$ such that $x\geq y$ and $x\reach y$,
	let $A'$ be 
the $d$-dim net obtained from $A$ by adding actions
	$\mathbf{e}_i$ for those $i\in[1,d]$ for which $x(i)>y(i)$, where $\mathbf{e}_i$
increments the
	$i$-th component by $1$.
	Formally,
	$(\mathbf{e}_i)_-=\mathbf{0}$ and $(\mathbf{e}_i)_+=e_i$,
	where $e_i$
is the $i$th unit vector satisfying
	$e_i(i)=1$ and $e_i(j)=0$ for 
	$j\neq i$.

Hence, there exist sequences $\sigma_1,\sigma_2$ of actions in $A'$
	satisfying
	\begin{equation}\label{eq:xmry}
	x\step{\sigma_1}y\step{\sigma_2}x\text{ in }A';
	\end{equation}
we can take $\sigma_1$ from $A$ due to the assumption $x\reach y$ in $A$,
whereas $\sigma_2$ can consist solely of the increment actions $\mathbf{e}_i$.
	By Lemma~\ref{lem:fmr}, there exist sequences
	$\sigma_1,\sigma_2$  of actions in $A'$
	that satisfy~\eqref{eq:xmry} and, moreover,
	\[|\sigma_1\sigma_2|\leq C\mr^A\cdot
\nnorm{x-y};\]
note that the actions $\mathbf{e}_i$ might occur
	in $\sigma_1$ as well. We safely use $C\mr^A=f\mr(\nnorm{A},d)$ instead of
	$C\mr^{A'}=f\mr(\nnorm{A'},d)$, since $\nnorm{A'}=\nnorm{A}$
	provided $\nnorm{A}\geq 1$; the claim is trivial if
	$\nnorm{A}=0$.

Any execution $x\step{\sigma_1}y$ in $A'$ can be clearly reorganized
	into the form 
	\[x\step{\sigma'}x'\step{\sigma}y,\] 
where 
the prefix $\sigma'$ is the projection of $\sigma_1$ onto
	the  set 
of increment actions $\mathbf{e}_i$, and
the suffix $\sigma$ is the projection of $\sigma_1$ onto
	the   set of original 
 actions in $A$.

	We have
	thus obtained $x'\geq x$ satisfying:
	\begin{enumerate}[a)]	
		\item $\nnorm{x'-x}\leq |\sigma'|\leq |\sigma_1|\leq C\mr^A\cdot \nnorm{x-y}$; 
		\item
			$x'\step{\sigma}y$ for some sequence $\sigma$ of actions in $A$
			such that $|\sigma|\leq |\sigma_1|\leq C\mr^A\cdot
			\nnorm{x-y}$;
	\item $x'(i)=x(i)$ if $x(i)=y(i)$ (since the action
		$\mathbf{e}_i$ is not in $A'$ for such $i\in[1,d]$).
			\qedhere
	\end{enumerate}
\end{proof}

\begin{lem}[\textbf{Large nonlive markings (virtually) reach  quasi-dead
	markings}]\label{lem:reachquasi}
	There exists a 2-exp function
$f\qd$ such that 
for every reversible $d$-dim net $A$ and 
	every nonlive marking
$x_0\in\upC{\max\{C\suf^A,\nnorm{A}\cdot C\qd^A\}}$, where
$C\qd^A=f\qd(\nnorm{A},d)$,
	there exists an $f\qd$-quasi-dead marking $x\in\Reach(x_0)$.
\end{lem}
\begin{proof}
Let $A$ be a reversible $d$-dim net and 
let	$x_0\in\upC{C\suf^A}$ be
	a nonlive marking of $A$.
	See Figure~\ref{fig:markingstwo} for illustration.
	The value $C\qd^A=f\qd(\nnorm{A},d)$ will arise from the
	following analysis.

	We fix a dead marking
	$y\in\Reach(x_0)$---it exists, since $x_0$ is nonlive---and a nonempty index set $I\subseteq [1,d]$ such that the
	restricted marking 
	$y\restr{I}\in[0,C\dead^A-1]^I$ is dead in the restricted net
	$A\restr{I}$ (recall~(\ref{eq:alldead})).
We have 
	\begin{equation}\label{eq:Iless}
		y(i)< x_0(i)\text{ for every }i\in I,
	\end{equation}		
because $y(i)< C\dead^A\leq C\suf^A\leq x_0(i)$ for every $i\in I$.

	Since  $x_0\reach y$, we obtain $y\vreach x_0$ by
the reversibility of $A$. We then consider a segmented
 virtual execution  
	\begin{equation}\label{eq:downup}
y=y_0\vreach y_1\vreach y_2\cdots\vreach y_k=x_0
	\end{equation}
	as in~(\ref{eq:virtsegment}).
	Hence, $\nnorm{y_{j+1}-y_j}\leq C\shift^A$ by~\eqref{eq:Cshift}, and 
for each $i\in[1,d]$ the following holds:
\begin{itemize}
\item If $y(i)<x_0(i)$, then
	$y(i)<y_1(i)<y_2(i)\cdots <y_{j_i}(i)=y_{j_i+1}(i)\cdots
	=y_k(i)=x_0(i)$;
\item	
If $y(i)\geq x_0(i)$, then $y_j(i)\geq x_0(i)$ for all
	$j\in[0,k]$.
	\end{itemize}
We have thus highlighted the strict increase in the case
	$y(i)<x_0(i)$, and the maintenance of the threshold $x_0(i)$ in
	the case $y(i)\geq x_0(i)$.

	Since $x_0(i)\geq C\suf^A$ for all $i\in[1,d]$, 
	there exists 	$j\leq C\suf^A$ such that 
	the marking $y_j$ in~\eqref{eq:downup}
	is above $C\suf^A$. By setting $x=y_j$, we can rewrite~(\ref{eq:downup}) as
	\begin{equation}\label{eq:xybound}
		y\vreach x\vreach x_0,\text{ where
		}x\in\upC{C\suf^A}\text{ and }\nnorm{x-y}\leq
		C\suf^A\cdot C\shift^A.
	\end{equation}
	We have thus obtained  a candidate for a quasi-dead marking $x\in\upC{C\suf^A}$, since 
	\begin{equation}\label{eq:candqd}
		x_0\reach x\reach y.
	\end{equation}
	Indeed, $x_0\reach y\vreach x$
	implies $x_0\vreach x$, hence  $x_0\reach x$ by
	Lemma~\ref{lem:fvr}, since 
	$x_0,x\in\upC{C\suf^A}$ and $C\suf^A=\max\{C\dead^A,C\vr^A\}\geq C\vr^A$. Due to
	reversibility, we also have $x\vreach x_0$, and thus  $x\reach
	x_0$, since 	$x,x_0\in\upC{C\vr^A}$.
Consequently, $x\reach y$ (because $x\reach x_0\reach y$).	
\begin{figure}[t]	
\begin{center}
\begin{tikzpicture}
    \begin{scope}
        \draw[black, thick] (0,0) -- (0,7);
        \draw[black, thick] (0,0) -- (8,0);
        \draw[black, thick] (8,0) -- (8,7);
        \node[anchor=south,rotate=90] at (0,3.5) {Number of tokens};

        \draw[dashed, black] (1.98,0) -- (1.98,7);
        \node[anchor=north] at (1,0) {$I$};
        \draw[dashed, black] (4.78,0) -- (4.78,7);
        \node[anchor=north] at (3.4,0) {$I' \smallsetminus I$};
        \node[anchor=north] at (6.5,0) {$[1,d] \smallsetminus I'$};
        
        \draw[densely dotted, black] (0,1.25) -- (8,1.25);
        \node[anchor=west] at (8,1.25) {$C\suf^A$};
        \draw[densely dotted, black] (0,4.5) -- (8,4.5);
	    \node[anchor=west] at (8,4.5) {$\nnorm{A}\cdot C\qd^A$};

        \drawconf{blue}{very thick}{
            6.9,4.5,6.1,5.5,6.5,
            6.0,6.1,5.7,5.0,4.8,
            5.0,5.4,5.9,6.22,6.16,
            4.9,4.5,4.7,5.1,5.0
        }{0.4}{0}
	    \node[anchor=west] at (8,5.0) {{\color{blue}$\bm{x_0}$} nonlive};

        \drawconf{applegreen}{very thick}{
            0.6,0.4,0.8,0.0,0.6,
            1.7,5.0,5.2,3.2,3.9,
            4.2,3.7,6.9,6.3,6.7,
            6.8,6.1,6.3,6.6,6.4
        }{0.4}{0}
	    \node[anchor=west] at (8,6.4)
	    {{\color{applegreen}$\bm{y}$} dead};

        \drawconf{red}{very thick}{
            1.9,1.5,1.8,1.7,1.3,
            2.3,5.8,5.65,3.95,4.6,
            4.95,4.3,6.35,6.26,6.2,
            6.2,5.2,5.9,5.8,6.0
        }{0.4}{0}
	    \node[anchor=west] at (8,6.0) {{\color{red}$\bm{x}$}
	    quasi-dead};
    \end{scope}
\end{tikzpicture}
\end{center}
	
\caption{Illustration for the proof of
	Lemma~\ref{lem:reachquasi}.}\label{fig:markingstwo}
\end{figure}
	

% Fix Indentation
\noindent
	However,
	we have not bounded the length of an execution from $x$ to a
	dead marking.
	To handle this, we recall that our choice of 
	the virtual execution $y\vreach{x}$ in~\eqref{eq:xybound} obtained
	from~(\ref{eq:downup}) guarantees that
	for each $i\in[1,d]$ we have
\begin{center}
either $y(i)< x(i)\leq x_0(i)$ or $y(i)\geq x(i)\geq x_0(i)$. 
\end{center}
	As illustrated in Figure~\ref{fig:markingstwo},
we consider the set
\[I'=\{i\in[1,d]\mid y(i)< x(i)\}=\{i\in[1,d]\mid y(i)<
	x_0(i)\}\supseteq I;\] 
	the inclusion $I\subseteq I'$ follows from~\eqref{eq:Iless}.
	 Since
	$y\restr{I}$ is dead in $A\restr{I}$, the marking 
$y\restr{I'}$ is dead in the net $A\restr{I'}$ (by
Proposition~\ref{prop:restrdead}). 

Because $x\restr{I'}\geq y\restr{I'}$ and  $x\restr{I'}\reach
y\restr{I'}$ in  $A\restr{I'}$ (due to $x\reach y$ in  $A$),
we can apply Proposition~\ref{prop:simpletrick} to the net
$A\restr{I'}$. We obtain  $x'\geq
x\restr{I'}$ (where $x'\in\setN^{I'}$) satisfying:
\begin{enumerate}[a)]
	\item 
$\nnorm{x'-x\restr{I'}}\leq 
		C\mr^{A\restr{I'}}\cdot
		\nnorm{x\restr{I'}-y\restr{I'}}\leq C\mr^{A}\cdot
		\nnorm{x-y}\leq  C\mr^{A}\cdot C\suf^A\cdot C\shift^A$
		(using~\eqref{eq:xybound}); and
		\item
			$x'\step{\sigma}y\restr{I'}$ for some $\sigma$ 
			such that $|\sigma|\leq C\mr^{A\restr{I'}}\cdot
	\nnorm{x\restr{I'}-y\restr{I'}}\leq C\mr^{A}\cdot C\suf^A\cdot C\shift^A$.
	\end{enumerate}
(Condition c) in Proposition~\ref{prop:simpletrick} is irrelevant
here.)

In $A\restr{I'}$, we thus have $x\restr{I'}\reach y\restr{I'}$
and $x'\reach y\restr{I'}$. 
By reversibility of $A\restr{I'}$ (inherited from $A$, recall
Proposition~\ref{prop:restrpresrev}), we obtain 
$x\restr{I'}\vreach x'$. Since $x'\geq x\restr{I'}\in\upC{C\suf^A}$
and $C\suf^A\geq
C\vr^A\geq C\vr^{A\restr{I'}}$,
Lemma~\ref{lem:fvr} yields  
 $x\restr{I'}\step{\sigma'} x'$ for a sequence $\sigma'$ of actions in
 $A\restr{I'}$, and thus in $A$ by our conventions, such that
\begin{equation}\label{eq:boundsigma}
	|\sigma'|\leq f\vr(\nnorm{A\restr{I'}},|I'|)\cdot
\nnorm{x'-x\restr{I'}}\leq C\vr^A\cdot
C\mr^A\cdot C\suf^A\cdot C\shift^A
\end{equation}
(using a)). To summarize, we have obtained the execution 
\begin{equation}\label{eq:finalseq}
	x\restr{I'}\step{\sigma'}x'\step{\sigma}y\restr{I'}\text{ in }A\restr{I'},
\end{equation}
where $y\restr{I'}$ is a dead marking, $|\sigma'|$ is bounded
in~\eqref{eq:boundsigma}, and  $|\sigma|$ is bounded
in~b).
Hence,
 $|\sigma'\sigma|\leq
f\qd(\nnorm{A},d)$, if we set
\begin{equation}\label{eq:fqd}
	f\qd(\nnorm{A},d)=C\qd^A=(C\vr^A+1)\cdot C\mr^A\cdot
C\suf^A\cdot C\shift^A.
\end{equation}
	Recall that $x(i)\geq x_0(i)$ for all $i\in[1,d]\smallsetminus
	I'$. 
	Consequently, if the considered nonlive marking 
\begin{center}	$x_0$ is above the level
	$\nnorm{A}\cdot C\qd^A$,
\end{center}
	we have 
$x(i)\geq\nnorm{A}\cdot C\qd^A$ for all $i\in[1,d]\smallsetminus I'$. 
Since $|\sigma'\sigma|\leq C\qd^A$, the sequence $\sigma'\sigma$ can be
executed not only from $x\restr{I'}$ in $A\restr{I'}$
(recall~\eqref{eq:finalseq}) but
	also from $x$ in $A$. We obtain an execution
	$x\step{\sigma'\sigma}y'$ in $A$, where $y'$ is a dead marking
	because ${y'}\restr{I'}=y\restr{I'}$ (and thus
	${y'}\restr{I}=y\restr{I}$, where $y\restr{I}$ is a dead
	marking in $A\restr{I}$).
The candidate $x$ in~\eqref{eq:candqd} is thus indeed an $f\qd$-quasi-dead marking, for the
2-exp function $f\qd$ defined by~\eqref{eq:fqd}.
\end{proof}	

\paragraph{Live solutions of $\bm{(A,x_0,\lambda)}$-linear systems.}
Further, we fix a 2-exp function $f\qd$ as guaranteed by
Lemma~\ref{lem:reachquasi}.  This fixes the constant
$C\qd^A=f\qd(\nnorm{A},d)$, and thus the notion of quasi-dead marking,
for every $d$-dim net $A$.
The next lemma refers to Definition~\ref{def:crucialLS}.

\begin{lem}[\textbf{$\bm{(A,x_0,\lambda)}$-solutions are live, for
	sufficiently large live $\bm{x_0}$ 
	and $\bm{\lambda}$}]\label{lem:finalbasicnets}
\hfill\\
	Let $A$ be a reversible $d$-dim net, and let
	$\lambda=(\lambda_1,\lambda_2,\dots,\lambda_d)$ be an extractor
	satisfying
	\begin{itemize}
		\item
			$\lambda_1\geq
C\suf^A+\nnorm{A}\cdot C\qd^A$;
\item
	$\lambda_{i+1}\geq \lambda_1 + \nnorm{A}\cdot C\vr^A\cdot
			(\lambda_i-C\suf^A)$ for each
			$i\in[1,d{-}1]$.
	\end{itemize}
Let $x_0$ be a live marking of $A$ satisfying	
	\[x_0\in\upC{\max\{C\suf^A,\nnorm{A}\cdot
	C\qd^A\}}.\]
Let $S$ be an $(A,x_0,\lambda)$-linear system, and let
	$(\bar{x}_0,(\bar{x}_I)_{\emptyset\neq I\subseteq[1,d]})$ be a
	solution of $S$.  Then $\bar{x}_0$ is a live marking.
\end{lem}	
\begin{proof}
	Let $A,\lambda$, and $x_0$ be as in the statement.
	Let $S$ be an $(A,x_0,\lambda)$-linear
	system~\eqref{eq:AxLS}, namely
\begin{equation*}
	(\bvar{x}\equiv_B x_0)
	\land
	\bigwedge_{\emptyset\neq I\subseteq[1,d]}\big((\bvar{x}_I\equiv_B x^\close_I)\land
	(\bvar{x}\virtual{*}\bvar{x}_I)\big),
\end{equation*}
with $B=\max\{C\suf^A,\lambda_d+C\shift^A\}$.
Recall that $({x}_0,({x}_I^\close)_{\emptyset\neq I\subseteq[1,d]})$ is a
solution of $S$, and that it is a tuple of live markings above
$C\suf^A$ (they are above $C\suf^A$ by
Definition~\ref{def:crucialLS}, and live by
Proposition~\ref{prop:liveunderreach}, since $x_0$ is live and
${x}_I^\close\in\Reach(x_0)$).
Towards a contradiction, suppose that
	\begin{center}
$(\bar{x}_0,(\bar{x}_I)_{\emptyset\neq I\subseteq[1,d]})$ is a solution of $S$
where $\bar{x}_0$ is nonlive. 
	\end{center}
	Note that for each $\emptyset\neq I\subseteq[1,d]$, we have
	$\bar{x}_0\vreach \bar{x}_I$ and
	$\bar{x}_I\equiv_B
	x^\close_I$, that is, 
	\begin{itemize}
		\item			
			$\bar{x}_I(i)=x^\close_I(i)$ \ for all
			$i\in[1,d]$ such that $x^\close_I(i)\leq B$, and
\item
			$\bar{x}_I(i)>B$ \ for all
			$i\in[1,d]$ such that $x^\close_I(i)>B$.
	\end{itemize}			
Since $x^\close_I\in\upC{C\suf^A}$ by definition, it follows that 
	$\bar{x}_I\in\upC{C\suf^A}$ as well.
	Moreover, the sets of indices of small components 
	$I^\lambda_{\bar{x}_I}$ and $I^\lambda_{x^\close_I}$
	(as defined in	Proposition~\ref{prop:Ilambdax}) are
	identical,
	and the markings $\bar{x}_I$ and $x^\close_I$ coincide on
	these indices.

\smallskip

	Lemma~\ref{lem:reachquasi} implies the existence of
	\begin{center}
	a quasi-dead
		marking $x\in\upC{C\suf^A}$ 
		such that 
	$\bar{x}_0\reach x$; let
 $I=I^\lambda_{x}$. 
	\end{center}		
Since $x$ is a quasi-dead marking, we have $x\step{\sigma_1}y$ where
$y$ is a dead marking and
	$|\sigma_1|\leq C\qd^A$.  
	Note that $I\neq\emptyset$. Indeed, we have $y(i)< C\dead^A$
	for at least one $i\in[1,d]$ (recall~\eqref{eq:alldead}).
The case $I=I^\lambda_{x}=\emptyset$ would imply,	 for all
	$i\in[1,d]$,
	\[x(i)\geq
	\lambda_1\geq C\suf^A+\nnorm{A}\cdot
	C\qd^A\geq 
	C\dead^A+\nnorm{A}\cdot |\sigma_1|;\]
this would in turn imply $y\in\upC{C\dead^A}$---a contradiction.  
Consequently, $\bar{x}_{I}$ is well-defined.

	Since all markings $\bar{x}_0$, $\bar{x}_I, x$ are above $C\suf^A$
and $A$ is
	reversible, we have 
	\[\bar{x}_0\reach \bar{x}_I\reach \bar{x}_0\reach
	x\reach\bar{x}_0\text{, hence }\bar{x}_I\reach x\text{ in
	particular.}\]
Note that $\dist{\lambda}{I}{\bar{x}_I}=\dist{\lambda}{I}{x^\close_I}=0$
(the distance to the $\lambda$-case $I$;
recall Figure~\ref{fig:extractor}(right)).
	The equality
	$\dist{\lambda}{I}{\bar{x}_I}=\dist{\lambda}{I}{x^\close_I}$
holds due to $\bar{x}_I\equiv_{B}x^\close_I$. If we had 
	$\dist{\lambda}{I}{\bar{x}_I}>0$, then 
	the first small segment $\bar{x}_I\virtual{\sigma}x'$
	of the virtual execution
	$\bar{x}_I\vreach x$ in the form
	of~(\ref{eq:virtsegment})---for which
	$\nnorm{x'-\bar{x}_I}\leq C\shift^A$ and 
	$x'(i)$ is between the values $\bar{x}_I(i)$ and $x(i)$ for
	each
	$i\in[1,d]$---would decrease this distance 
(that is, $\dist{\lambda}{I}{x'}<\dist{\lambda}{I}{\bar{x}_I}$).
For the sequence $\sigma$ from $\bar{x}_I\virtual{\sigma}x'$, we obtain 
 $x^\close_I\virtual{\sigma}x''$, where  $x'\equiv_{\lambda_d}x''$
 (since  $\bar{x}_I\equiv_{B}x^\close_I$ and
$B=\lambda_d+C\shift^A$). Consequently, 
$x^\close_I\reach x''$ (as both $x^\close_I$ and $x''$ are above
$C\suf^A\geq C\vr^A$), 
and
$\dist{\lambda}{I}{x''}<\dist{\lambda}{I}{x^\close_I}$; this contradicts the choice
of  $x^\close_I$.
We thus have
\[\dist{\lambda}{I}{x}=\dist{\lambda}{I}{\bar{x}_I}=\dist{\lambda}{I}{x^\close_I}=0,\]
and all markings $x,\bar{x}_I,x^\close_I$ are above $C\suf^A$.
Hence, for each $i\in I$, the values $x(i)$ and
$\bar{x}_I(i)=x^\close_I(i)$ are in the interval
$[C\suf^A,\lambda_{|I|}-1]$.
The fact 
$\bar{x}_I\reach x$ implies that in the restricted net $A\restr{I}$ we
have 
$(\bar{x}_I)\restr{I}\step{\sigma_2}x\restr{I}$ where 
\[|\sigma_2|\leq f\vr(\nnorm{A},|I|)\cdot
\nnorm{x\restr{I}-(\bar{x}_I)\restr{I}}\leq C\vr^A\cdot
(\lambda_{|I|}-C\suf^A).\]
Combining with $x\step{\sigma_1}y$, where $|\sigma_1|\leq
C\qd^A$ and $y$ is dead, we obtain
\[(\bar{x}_{I})\restr{I}=(x^\close_{I})\restr{I}\step{\sigma_2}
x\restr{I}\step{\sigma_1}y\restr{I}\]
in $A\restr{I}$, where $|\sigma_2\sigma_1|\leq C\vr^{A}\cdot
(\lambda_{|I|}-C\suf^A)+C\qd^A$.
Since for each $i\in [1,d]\smallsetminus I$, we have
$x^\close_{I}(i)\geq\lambda_{|I|+1}\geq \nnorm{A}\cdot
|\sigma_2\sigma_1|$, it follows that 
\[x^\close_{I}\step{\sigma_2\sigma_1}y'\]
is an execution in $A$, where $y'\restr{I}=y\restr{I}$.
We complete the proof by showing that $y'$ is dead, which
contradicts the assumption that $x^\close_{I}$ is live (since 
$x_0\reach x^\close_{I}$ and $x_0$ is live).

Recall that $x\step{\sigma_1}y$ and $|\sigma_1|\leq C\qd^A$.
Since for each $i\in [1,d]\smallsetminus I$, we have 
\[x(i)\geq
\lambda_{|I|+1}\geq \lambda_1\geq C\dead^A+ \nnorm{A}\cdot C\qd^A,\]
it follows that 
$y(i)\geq C\dead^A$ for each  $i\in [1,d]\smallsetminus I$. This implies
that $y\restr{I}$ is dead in $A\restr{I}$, which in turn
implies that $y'$ is dead in $A$ (since
$y'\restr{I}=y\restr{I}$).
\end{proof}

\paragraph{Petri nets with states (PNSs).}
Lemma~\ref{lem:finalbasicnets}, which gives conditions under which 
solutions of $(A,x_0,\lambda)$-linear systems are live,
and Proposition~\ref{prop:Affsmall}, which guarantees 2-exp solutions of 
such systems, together establish Theorem~\ref{th:upper} 
for reversible nets $A$ that have a live marking $x_0$ above the
level $\max\{C\suf^A, \nnorm{A}\cdot C\qd^A\}$.

However, there exist structurally live conservative nets 
(hence, nets with live bottom SCCs)
that have no live marking above that level.
Indeed, there are nets in which the values of some 
components---indexed by a~set $\emptyset\neq I\subseteq[1,d]$---remain
small in every live marking, as demonstrated
by Example~\ref{exa:basiclive}.
In this case, it is natural to view the restrictions $x\restr{I}$ of
(live) markings $x$ as \emph{control states}; the components (that is,
Petri net places) outside $I$ are viewed as \emph{counters}.
Definition~\ref{def:PNS}
introduces \emph{Petri nets with
states} (PNSs), in a specific form tailored to our use.   
We will extend the previous results from nets to PNSs, with live
markings having sufficiently large counter values.

\begin{rem}
Petri nets with states can be viewed, in principle, 
as \emph{vector addition
	systems with states} (VASSs); a minor difference is 
	that transitions between states are
labelled with Petri net actions $a=(a_-,a_+)$ in PNSs, whereas in
	VASSs they are
	labelled merely with 
plain vectors encoding the effects $\Delta(a)$. 
	Our
	definition of PNSs is also related to the notion of
	\emph{rigid components} introduced in~\cite{DBLP:conf/stoc/Kosaraju82} in order to keep track 
	of the exact values of 
bounded components.
\end{rem}

% Fix Indentation
\noindent
Given a $d$-dim net $A$ and $I\subseteq[1,d]$,
we say that a \emph{nonempty finite} set $Q\subseteq\setN^I$ is
a~\emph{bSCC} of the restricted net $A\restr{I}$
if it corresponds to a bottom SCC 
in the
reachability graph of $A\restr{I}$; that is, if in $A\restr{I}$ we have 
$\Reach(x)=Q$ for all $x\in Q$.
Figure~\ref{fig:example} shows a 5-dim net $A_1$ and a restricted 4-dim
 net
$A_2=(A_1)\restr{[1,4]}$. Recalling Example~\ref{exa:basiclive}, we
can verify that $Q=\{(1,0,1,0),(0,1,2,0),(0,1,0,1),(1,0,0,1)\}$ is a
bSCC of $A_2$.

\begin{defi}[\textbf{PNS $\bm{G^A_{(I,Q)}}$, norm $\bm{\nnorm{G}}$,
	relations $\bm{\reach}$ and $\bm{\vreachtitle}$,
	reversible PNS}]\label{def:PNS}
	Given a $d$-dim net $A$, an index set $I\subseteq [1,d]$, and a (finite) bSCC 
$Q\subseteq\setN^I$, 
the \emph{Petri net with states (PNS)} $G^A_{(I,Q)}$ 
	is obtained from  $A$ by restricting \emph{markings}
	$\bar{x}\in\setN^d$ 
	to those satisfying $\bar{x}\restr{I}\in Q$. 
Setting  $J=[1,d]\smallsetminus I$,
	such a~marking  $\bar{x}\in\setN^d$
	is also written as 
	\begin{equation*}
		\textnormal{$(p,x)\in
	Q\times \setN^J$ where 
		$p=\bar{x}\restr{I}\in Q$ and
		$x=\bar{x}\restr{J}\in\setN^J$.}
	\end{equation*}
For $G=G^A_{(I,Q)}$, we say that $G$ is of \emph{dimension} $d$
(inherited from $A$), or that $G$ is a~\emph{$d$-dim PNS}.
We define the \emph{norm of} $G$ as
	\[\nnorm{G}=\max\{\nnorm{Q},\nnorm{A}\}.\]
The relations $\step{a}$,
	$\step{\sigma}$, $\reach$ in $G^A_{(I,Q)}$  coincide with the respective
	relations in $A$ restricted to markings in $Q\times \setN^J$.
	The \emph{virtual reachability relation} in  $G^A_{(I,Q)}$ is defined 
on the set $Q\times \setZ^J$ as follows:
	for action $a$ in $A$ we have 
	\begin{center}
	$(p,x)\virtual{a}(q,y)$ if $p\step{a\restr{I}}q$ in
	$A\restr{I}$ and $x\virtual{a\restr{J}}y$ in $A\restr{J}$.
	\end{center}
	This naturally induces the
	relations $(p,x)\virtual{\sigma}(q,y)$ for sequences
	$\sigma$ of
	actions in $A$, and the virtual reachability relation
	$(p,x)\vreach(q,y)$ on the set $Q\times \setZ^J$. 
Virtual reachability in a PNS $G^A_{(I,Q)}$ is thus based on
the standard reachability within $Q$ in $A\restr{I}$:
	\begin{center}
	$(p,x)\virtual{\sigma}(q,y)$ if and only if $p\step{\sigma}q$ in
		$A\restr{I}$ and $y-x =\Delta(\sigma)\restr{J}$.
	\end{center}
(By our
	convention, we may write 
	 $p\step{\sigma}q$ as a shorthand for
	  $p\step{\sigma\restr{I}}q$.)



	A \emph{PNS} $G^A_{(I,Q)}$ is \emph{reversible} if its virtual
	reachability relation
	$\vreach$ is symmetric; that is, $(p,x)\vreach(q,y)$ implies 
	$(q,y)\vreach(p,x)$.
\end{defi}


\begin{rem}
The dimension of $G=G^A_{(I,Q)}$, where $A$ is a $d$-dim net and
	$I\subseteq [1,d]$,
	could be naturally defined as $|J|$ for $J=[1,d]\smallsetminus
	I$. This would correspond to
the standard notion 
	of dimension in vector addition systems with states (VASSs). 
However, we define it as $d$ to emphasize that the underlying net $A$
remains the primary object of study.  

Note that the requirement that the state set $Q$ in $G=G^A_{(I,Q)}$ is a bSCC of
	$A\restr{I}$ guarantees that for any marking
	$\bar{x}=(p,x)\in Q\times \setN^J$, the reachability set $\Reach(\bar{x})$ in $A$ coincides with that in $G$.
\end{rem}

\paragraph{Extractors $\bm{\lambda}$ yielding reversible PNSs
$\bm{G^{(A,\lambda)}_{(I,Q)}}$ for
live bottom SCCs of nets $\bm{A}$.}

The following lemma is a simplified variant of Lemma~20 in
Section~6 of~\cite{DBLP:conf/fsttcs/Leroux19}. It concerns a general
SCC of a net; a bottom SCC 
is a special case. It allows us to separate, within an SCC, the
``control states''---components whose values are small in every marking
of the SCC---from the ``counters'', whose values are simultaneously
large in at least one marking.
The lemma uses the notion of \emph{$\nnorm{A}$-adapted
extractors}. For later convenience, we also introduce the notion of
\emph{strictly $\nnorm{A}$-adapted
extractors}.

\begin{defi}[\textbf{$\bm{\nnorm{A}}$-adapted and strictly $\bm{\nnorm{A}}$-adapted extractors
	$\bm{\lambda}$}]\label{def:Aadaptextract}
\hfill\\
	Given a $d$-dim extractor
	$\lambda=(\lambda_1,\lambda_2,\dots,\lambda_d)$, with
	$\lambda_0=1$ by our convention,
	and
	$m\in\setN_+$ (in particular, $m=\nnorm{A}$ for a $d$-dim net
	$A$), we say that $\lambda$ is \emph{$m$-adapted}
if  for each
	$i\in[0,d{-}1]$,
	\[\lambda_{i+1}\geq \lambda_i+m\cdot (\lambda_i)^i.\]
	 We say that $\lambda$ is  \emph{strictly
	$m$-adapted} if $\lambda_1\geq 1+m$ and
	for each
	$i\in[1,d{-}1]$,
	\begin{equation}\label{eq:genextr}
		\lambda_{i+1}\geq m\cdot (\lambda_i)^i+m\cdot
	i\cdot(\lambda_i)^i.
	\end{equation}
	(Note that every strictly $m$-adapted extractor is $m$-adapted.)
\end{defi}


\begin{lemC}[\cite{DBLP:conf/fsttcs/Leroux19} (\textbf{Small control-state projection with large
	counters})]\label{lem:sccfirst}
\hfill\\
	Let $X$ be an SCC in the reachability graph of a $d$-dim net $A$ 
 and
	let
	$\lambda=(\lambda_1,\lambda_2,\dots,\lambda_d)$ be a~$d$-dim
	extractor that is $\nnorm{A}$-adapted. 
	There exists a 
	set $I\subseteq[1,d]$
	satisfying the following conditions, where
	$J=[1,d]\smallsetminus I$:
	\begin{enumerate}[a)]	
\item
	$\nnorm{\rst{X}{I}}<\lambda_{|I|}$ (hence for all $x\in X$ and
			$i\in I$ we have $x(i)<\lambda_{|I|}$);			
\item
there exists			$x\in X$
such that for all $j\in J$ we have
		$x(j)\geq
		\lambda_{\sizeof{I}+1}-\nnorm{A}\cdot\sizeof{I}\cdot(\lambda_{\sizeof{I}})^{\sizeof{I}}$.
	\end{enumerate}
\end{lemC}


% Fix Indentation
\noindent
In the previous lemma, the mere existence of a respective set $I\subset[1,d]$
is sufficient for our aims; we do not require its uniqueness in
this case. This entails that we do not require the uniqueness of 
the PNS $G^{(A,\lambda)}_{(I,Q)}$ in the following definition.

\begin{defi}[\textbf{Extracting PNS
	$\bm{G^{(A,\lambda)}_{(I,Q)}}$ from a live bottom SCC of
	a net $\bm{A}$}]\label{def:extrPNS}
\hfill\\
Let a (finite or infinite) set $X\subseteq \setN^d$ be a live bottom SCC in the reachability graph of
	a $d$-dim net $A$. (The net $A$ is reversible by
	Lemma~\ref{lem:consstrbound}.)
	Let $\lambda$ be a strictly $\nnorm{A}$-adapted $d$-dim
extractor, and let $I\subseteq[1,d]$ be a set guaranteed by
Lemma~\ref{lem:sccfirst}. 
We then view the PNS $G=G^A_{(I,Q)}$ with  $Q=X\restr{I}$ as a PNS 
	$G^{(A,\lambda)}_{(I,Q)}$, \emph{extracted from the SCC $X$ via the extractor
	$\lambda$}.
\end{defi}


% Fix Indentation
\noindent
In the previous definition, we consider strictly $\nnorm{A}$-adapted 
extractors to guarantee that $Q=X\restr{I}$ is a bSCC of $A\restr{I}$
(as required in Definition~\ref{def:PNS}):

\begin{prop}[\textbf{PNSs $\bm{G^{(A,\lambda)}_{(I,Q)}}$ are
	well-defined, and reversible}]\label{prop:condproperPNS}
\hfill \\	
In every PNS ${G^{(A,\lambda)}_{(I,Q)}}$,
	the set ${Q}$ is indeed a
	bSCC of ${A\restr{I}}$. Moreover, ${G^{(A,\lambda)}_{(I,Q)}}$
	is reversible.
\end{prop}	
\begin{proof}
Let a PNS $G={G^{(A,\lambda)}_{(I,Q)}}$ be extracted from a live bottom
	SCC $X$ of a $d$-dim net $A$ as described
	in Definition~\ref{def:extrPNS}; this implies
that $A$ is reversible,	$Q=X\restr{I}$, and $\nnorm{Q}<\lambda_{|I|}$ (by
	condition a) in Lemma~\ref{lem:sccfirst}). Consequently, 
	\[|Q|\leq (1+\nnorm{Q})^{|I|}\leq (\lambda_{|I|})^{|I|}.\]

	If $I=\emptyset$, then $Q=\{()\}$ (the empty tuple is the
	unique
	control state), and we have $()\step{a}()$ for all actions in
	$A$. Hence, $Q$ is a trivial bSCC of $A\restr{I}$, and $G$ is
	reversible because $A$ is reversible.
	If $I=[1,d]$, then $A\restr{I}=A$ and $Q=X$.
	Here, the facts that $Q$ is a bSCC of $A\restr{I}$ and that $G$ is reversible follow from $X$ being a bottom SCC of $A$.
	We thus further assume $\emptyset\neq I\neq[1,d]$.
	
	To prove that  ${Q}=X\restr{I}$ is a bSCC of ${A\restr{I}}$,
we assume $p\step{a}p'$ 
	(that is, $p\step{a\restr{I}}p'$) 
	in $A\restr{I}$ with
	$p\in Q$, and show that $p'\in Q$; more precisely, we show
	that there exists a step 
	$(p,y)\step{a}(p',y')$ within $X$ for some
	$y,y'\in\setN^J$, where $J=[1,d]\smallsetminus I$.
	We use the fact that by condition b) in
	Lemma~\ref{lem:sccfirst}, there exists 
	$(q_0,x_0)\in X$ such that ($q_0\in Q$ and) for each $j\in J$,
	\[
	x_0(j)\geq
		\lambda_{\sizeof{I}+1}-\nnorm{A}\cdot\sizeof{I}\cdot(\lambda_{\sizeof{I}})^{\sizeof{I}}\geq
	\nnorm{A}\cdot (\lambda_{|I|})^{|I|}\geq \nnorm{A}\cdot |Q|.\]
(We use the property~\eqref{eq:genextr} of strict
	$\nnorm{A}$-adapted extractors.)
	Since $Q=X\restr{I}$, $X$ is an SCC, and $q,p\in Q$, there exists an execution 
	\[q_0\step{a_1}q_1\step{a_2}q_2\cdots \step{a_k}q_k=p\] in $A\restr{I}$ 
	with $k\leq |Q|-1$. The corresponding execution
	\[(q_0,x_0)\step{a_1}(q_1,x_1)\step{a_2}(q_2,x_2)\cdots
	\step{a_k}(q_k,x_k)=(p,y)\]
	 is contained within $X$, because $x_0$
	is above $\nnorm{A}\cdot |Q|\geq (k+1)\cdot\nnorm{A}$, and
	thus each $x_i$ is above
	$(k+1-i)\cdot\nnorm{A}\geq\nnorm{A}$.
The step
	$(p,y)\step{a}(p',y+\Delta(a)\restr{J})$ is thus within $X$ as
	well. 

	The proved fact that $p\step{a}p'$ with $p\in Q$ implies  a step $(p,y)\step{a}(p',y')$
	within $X$ shows not only 
	that $Q$ is a bSCC of $A\restr{I}$, but also  
that	${G^{(A,\lambda)}_{(I,Q)}}$ is reversible: 
	If
	$(p,x)\virtual{a}(p',x')$, that is, $p\step{a\restr{I}}p'$ and
	$x'-x=\Delta(a)\restr{J}$, then within the SCC $X$ we have
	$(p,y)\step{a}(p',y')$
	with $y'-y=x'-x=\Delta(a)\restr{J}$; this implies the
	existence of an execution
	$(p',y')\step{\sigma} (p,y)$ (within $X$), and thus
	$(p',x')\virtual{\sigma} (p,x)$.
\end{proof}	

\paragraph{Adapting the proof of Theorem~\ref{th:upper} from 
reversible nets to reversible PNSs}
The issue that a $d$-dim net $A$ with a live bottom SCC $X$ might not
possess a live marking 
above a sufficiently large level---specifically above
$\max\{C\suf^A,\nnorm{A}\cdot C\qd^A\}$ as stated in
Lemma~\ref{lem:finalbasicnets}---is, in principle,
resolved by extracting a PNS $G=G_{(I,Q)}^{(A,\lambda)}$ from $X$.
For a suitably chosen 2-exp extractor $\lambda$, the norm
$\nnorm{G}=\max\{\nnorm{Q},\nnorm{A}\}$ is (at most) 2-exp, and there
exists a live marking $(p_0,x_0)$ with $x_0$ above a sufficiently
large level (by part b) of Lemma~\ref{lem:sccfirst}).

In what follows, we therefore adapt the arguments from the previous proof
for reversible nets---which can be viewed as reversible PNSs with a
single state (the empty tuple)---to the full case of reversible PNSs.


\paragraph{(Virtual) reachability, 1-exp
$\bm{f\bvr}$ and $\bm{C\bvr^G}$, 2-exp
$\bm{C\bsuf^G=\max\{C\bdead^A,C\bvr^G\}}$.}
We now recall a variant of~\cite[Lemma~5]{DBLP:conf/fsttcs/Leroux19} in a more
general form than in Lemma~\ref{lem:fvr}:

\begin{lem}[\textbf{Virtual and standard reachability in
reversible PNSs, $\bm{f\bvr}$ and $\bm{C\bvr^G}$}]\label{lem:pnsreachfirst}
	There exists a 1-exp function $f\vr$ such that 
for every reversible $d$-dim PNS $G$, the following conditions hold,
	where	
$C\vr^G=f\vr(\nnorm{G},d)$:
	\begin{enumerate}
		\item
			If $(p,x)\vreach (q,y)$, then
			$(p,x)\virtual{\sigma}(q,y)$
for some $\sigma$ satisfying
			$|\sigma|\leq
					C\vr^G\cdot \nnorm{y-x}$.
\item
	If $(p,x)\vreach (q,y)$
			and  $x,y\in\upC{C\vr^G}$, 
			then $(p,x)\step{\sigma}(q,y)$
	with
		$|\sigma|\leq
		C\vr^G\cdot \nnorm{y-x}$.
	\end{enumerate}


\end{lem}
\begin{rem}
The cited result in~\cite{DBLP:conf/fsttcs/Leroux19} concerns
a broader class 
of reversible PNSs than those considered here and is technically more
detailed.
As in other parts of this paper, we prioritize a simplified presentation 
that suffices for our needs.
\end{rem}


% Fix Indentation
\noindent
Dead markings of a PNS $G=G_{(I,Q)}^A$ are the dead markings $y$ of the
net $A$ satisfying $y\restr{I}\in Q$. Therefore, we do not introduce
a separate constant
$C\dead^G$, but simply use $C\dead^A$.  The sufficiently large 2-exp
constant from~\eqref{eq:Csuf} now becomes
\begin{equation}\label{eq:CsufG}
C\suf^G=\max\{C\dead^A,C\vr^G\}.
\end{equation}

\paragraph{Virtual reachability in reversible PNSs via simple cycles
in control-state units.}
Control-state units of PNSs can be naturally viewed as graphs, allowing us
to use standard graph-theoretic terminology:

\begin{defi}[\textbf{Arc-labelled multigraph $\bm{(Q,E)}$ corresponding to
	a PNS $\bm{G_{(I,Q)}^A}$}]
	Given a $d$-dim PNS $G=G_{(I,Q)}^A$, by $(Q,E)$ we denote 
the arc-labelled multigraph 
over the set of actions $A$, 	where 
$p\step{a}q$ is an arc in $E$ if $p\step{a\restr{I}}q$ holds in
the net $A\restr{I}$.
\\
By a \emph{path} in $(Q,E)$, we mean a sequence
	$p_0\step{a_1}p_1\cdots \step{a_k}p_k$; it is a \emph{cycle}
	if $p_k=p_0$, and a \emph{simple
cycle} if, moreover, $0\leq i<j\leq k-1$ implies $p_i\neq p_j$. 
By the
\emph{effect} of a cycle $p\step{\sigma}p$, we mean
	the vector $\Delta(\sigma)\restr{J}\in\setZ^J$, where
	$J=[1,d]\smallsetminus I$
(noting that $\Delta(\sigma)\restr{I}=\mathbf{0}$).
\end{defi}	

% Fix Indentation
\noindent
The following proposition motivates  focusing on
 ``cyclic virtual reachability''.

\begin{prop}[\textbf{Virtual reachability in PNSs is based on
	control-unit cycles}]\label{prop:quickreachcycle}
Given a reversible $d$-dim PNS $G=G_{(I,Q)}^A$, if $(p,x)\vreach (q,y)$, then
there exists an action sequence $\sigma$ with $|\sigma|\leq
	(|Q|-1)$---which implies $\nnorm{\Delta(\sigma)}\leq
	(|Q|-1)\cdot\nnorm{A}$---such that
	\[(p,x)\virtual{\sigma}(q,x')\vreach(q,y),\]
where $x'=x+\Delta(\sigma)\restr{J}$ for $J=[1,d]\smallsetminus I$.
\end{prop}
\begin{proof}
If $(p,x)\vreach (q,y)$ in a reversible PNS $G=G_{(I,Q)}^A$,
we may simply consider a
	shortest path $p\step{\sigma}q$ in the corresponding graph $(Q,E)$;
we thus obtain $(p,x)\virtual{\sigma}(q,x')$ with
	$|\sigma|\leq |Q|-1$.
By reversibility of $G$, the relation $\vreach$ is symmetric, which
	implies  $(q,x')\vreach(p,x)$; together with the assumption $(p,x) \vreach
(q,y)$, this yields $(q,x') \vreach (q,y)$. 
\end{proof}


% Fix Indentation
\noindent
We note that simple cycles constitute the building blocks of
cyclic virtual executions:

\begin{prop}[\textbf{Zero-effect all-state cycles, and simple cycles as
	basic blocks}]\label{prop:cycles}
	Given a reversible $d$-dim PNS $G=G_{(I,Q)}^A$, the following
	conditions hold:
	\begin{enumerate}
		\item	There exists a cycle
	$p_0\step{\sigma_0}p_0$ in the corresponding graph $(Q,E)$
	that visits all states in $Q$ and has zero effect; that
			is, it satisfies 
			not only
	$\Delta(\sigma_0)\restr{I}=\mathbf{0}$ but also
	$\Delta(\sigma_0)\restr{J}=\mathbf{0}$ for
			$J=[1,d]\smallsetminus I$.
\item 
	For every state $p_0\in Q$, the effects of cycles
			$p_0\step{\sigma}p_0$
			are precisely those vectors in $\setZ^J$ that
			can be expressed as sums of the effects of
			simple cycles in $(Q,E)$.
	\end{enumerate}
	\end{prop}
\begin{proof}
	(1) By definition, $Q$ is a bottom SCC in $A\restr{I}$.
	Consequently, there is a cycle $p\step{\sigma}p$ in the graph
	$(Q,E)$ that visits all states $q\in Q$; the cycle corresponds to a virtual
	execution
	$(p,\mathbf{0})\virtual{\sigma}(p,\Delta(\sigma)\restr{J})$.
	By the symmetry of the relation $\vreach$, there must also
	exist a virtual execution
	$(p,\Delta(\sigma)\restr{J})\virtual{\sigma'}(p,\mathbf{0})$
	for some sequence $\sigma'$. Therefore,
	$p\step{\sigma\sigma'}p$ is the
	claimed zero-effect cycle visiting all states.

	(2) Let $\mathcal{E}\subseteq \setZ^J$ be the set
	of effects of simple cycles in the graph $(Q,E)$.
Any cycle $p_0\step{\sigma}p_0$ in 
	$(Q,E)$ is either a~simple cycle---in which case
	$\Delta(\sigma)\restr{J}\in\mathcal{E}$---or it can be
	decomposed as 
	$p_0\step{\sigma_1}q\step{\sigma_2}q\step{\sigma_3}p_0$, where
	$q\step{\sigma_2}q$ is a simple cycle. In the latter case, the
	effect of the cycle $p_0\step{\sigma_1\sigma_3}p_0$ belongs to
	$\mathcal{E}^*$ by the inductive hypothesis; since
	$\Delta(\sigma_2)\restr{J}\in\mathcal{E}$, it follows
	that $\Delta(\sigma_1\sigma_2\sigma_3)\restr{J}$ belongs to
	$\mathcal{E}^*$ as well. 

	On the other hand, if $z_0\in \mathcal{E}$---due to 
	 a simple
	cycle $q\step{\sigma}q$ with $\Delta(\sigma)\restr{J}=z_0$---
	then for any $p_0\in Q$, we may consider a cycle
	$p_0\step{\sigma_1}q\step{\sigma}q\step{\sigma_2}p_0$, where 
	$p_0\step{\sigma_1\sigma_2}p_0$ is a zero-effect cycle visiting
	all states in $Q$. Consequently, the effect of the cycle
	$p_0\step{\sigma_1\sigma\sigma_2}p_0$ is $z_0$.
	This implies that for any
	$z\in\mathcal{E}^*$, there exists a cycle 
	$p_0\step{\sigma}p_0$ with the effect $\Delta(\sigma)\restr{J}=z$.
\end{proof}


\paragraph{Cyclic virtual reachability in a PNS $\bm{G}$ via a
``simple-cycle'' net $\bm{\ASC^G}$.}
Recall that a $d$-dim net $A$ is a finite set of actions
 $a=(a_-,a_+)\in\setN^d\times\setN^d$, with
$\Delta(a)=(a_+-a_-)\in\setZ^d$. For the virtual reachability relation $\vreach$,
only the set $A\D=\{\Delta(a)\mid a\in A\}$ is relevant, since
$x\vreach y$ holds
if and only if $(y-x)\in A\D^*$. 
In the following definition, for $z\in \setZ^J$ we
use the notation $z_+$, $z_-$ where, for each $i\in J$,
\begin{center}
$z_+(i)=\max\{z(i),0\}$ and $z_-(i)=\max\{-z(i),0\}$.
\end{center}
	For instance, $(2,-3,0)_+=(2,0,0)$ and  $(2,-3,0)_-=(0,3,0)$.
Thus $z=z_+-z_-$, and each $z\in\setZ^J$ can be implicitly viewed as
a net action $z=(z_+,z_-)$ with $\Delta(z)=z$.


\begin{defi}[\textbf{Simple-cycle \bm{$|J|$}-dim net \bm{$\ASC^G$} associated with
	a PNS \bm{$G=G^A_{(I,Q)}$}}]\label{def:ASC}
	Given a $d$-dim PNS $G=G^A_{(I,Q)}$ with the corresponding
	graph $(Q,E)$
and $J=[1,d]\smallsetminus I$, 
	the \emph{simple-cycle net} $\ASC^G$
is the $|J|$-dim net 
	\begin{equation*}
		\ASC^G=(\ASC^G)\D=\{z\in\setZ^J\mid z \text{ is the effect of a simple
		cycle in }(Q,E)\}.
	\end{equation*}

\end{defi}
% Fix Indentation
\noindent
We note that, given a $d$-dim PNS $G=G^A_{(I,Q)}$, we have
\begin{equation}\label{eq:normASC}
\nnorm{\ASC^G}\leq |Q|\cdot\nnorm{A}\leq
	\nnorm{A}\cdot(1+\nnorm{Q})^{|I|}.
\end{equation}
The following fact is a corollary of
Proposition~\ref{prop:cycles}(2).

	\begin{cor}[\textbf{Reducing cyclic virtual reachability in
		$\bm{G}$ 
		to the net 
		$\bm{\ASC^G}$}]\label{cor:cycles}
\hfill\\
		Given a reversible $d$-dim PNS $G=G^A_{(I,Q)}$, for any $p\in Q$ we
		have 
		\begin{center}
		$(p,x)\vreach (p,y)$ in $G$ if and only if $x\vreach y$
		in $\ASC^G$.
		\end{center}			
	\end{cor}

\paragraph{Extending $\bm{(A,x_0,\lambda)}$-linear systems to PNSs.}
We now extend the notion of $(A,x_0,\lambda)$-linear systems
from Definition~\ref{def:crucialLS} to PNSs.

\begin{defi}[\textbf{$\bm{(G,(p_0,x_0),\lambda)}$-linear systems for
	reversible PNSs}]\label{def:crucialLSPNS}
\hfill\\
	Given a reversible $d$-dim PNS	$G=G^A_{(I,Q)}$, where
	$J=[1,d]\smallsetminus I$, and a $|J|$-dim extractor
	$\lambda=(\lambda_1,\lambda_2,\dots,\lambda_{|J|})$,
	we define
the bound	\begin{equation}\label{eq:BPNS}
	B=B_{(G,\lambda)}=\max\{C\suf^G,\lambda_{|J|}+C\shift^{\ASC^G}\}.
	\end{equation}
	Given a marking $(p_0,x_0)\in Q\times \setN^J$ 
	with $x_0\in\upC{C\suf^G}$,
	for each $\emptyset\neq J'\subseteq J$ let 
	\[x^\close_{J'}\in \setN^J\cap\upC{C\suf^G}\]
	be a fixed vector such that
	$(p_0,x_0)\reach(p_0,x^\close_{J'})$ and  
	$\dist{\lambda}{J'}{x^\close_{J'}}$ is
	minimal. 
Then the following 
	linear system $S$ of dimension $|J|\cdot 2^{|J|}$, with
	a $|J|$-dim variable
	$\bvar{x}=(\bvar{x}_1,\bvar{x}_2,\dots,\bvar{x}_{|J|})$ and
	$2^{|J|}-1$ $|J|$-dim variables 
	$\bvar{x}_{J'}$, one for each $\emptyset\neq J'\subseteq J$,
	is 
	\begin{center}
		a \emph{$(G,(p_0,x_0),\lambda)$-linear system}:
	\end{center}		
\begin{equation}\label{eq:GxLS}
	(\bvar{x}\equiv_B x_0)
	\land
	\bigwedge_{\emptyset\neq J'\subseteq
	J}\big((\bvar{x}_{J'}\equiv_B x^\close_{J'})\land
	(\bvar{x}\virtual{*}\bvar{x}_{J'})\big),
\end{equation}
	where $\bvar{x}\virtual{*}\bvar{x}_{J'}$ refers to the virtual
	reachability in $\ASC^G$. 
\end{defi}	


% Fix Indentation
\noindent
Note that $(x,x')$ is a solution of the expression
$\bvar{x}\virtual{*}\bvar{x}_{J'}$ referring to $\ASC^G$ if and
only if $(p_0,x)\vreach(p_0,x')$ holds in $G$ (due to Corollary~\ref{cor:cycles}).

Our aim is to show that if $(p_0,x_0)$ is a live marking of
$G^A_{(I,Q)}$, where $x_0$ is above a sufficiently large level
relative to $\nnorm{G}=\max\{\nnorm{Q},\nnorm{A}\}$, then  $(p_0,\bar{x}_0)$ is live for every
solution $(\bar{x}_0,(\bar{x}_{J'})_{\emptyset\neq J'\subseteq J})$
of a $(G,(p_0,x_0),\lambda)$-linear system, for a suitable 2-exp
extractor $\lambda$.
As previously, we use the notion of quasi-dead markings to this aim.


\begin{defi}[\textbf{Quasi-dead markings in PNSs}]
\hfill\\	
	For a 2-exp function $f\qd$ and a $d$-dim PNS $G=G^A_{(I,Q)}$,
	where $J=[1,d]\smallsetminus I$,
	a marking $(p_0,x)\in Q\times \setN^J$
	is \emph{$f\qd$-quasi-dead} in $G$, or simply
	\emph{quasi-dead} in $G$ when $f\qd$
	is clear from context, if $x\in\upC{C\suf^G}$ 
	and there exists an execution $(p_0,x)\step{\sigma}(q,y)$ such
	that $(q,y)$ is a dead
marking and $|\sigma|\leq
f\qd(\nnorm{G},d)=C\qd^G$.
\end{defi}

\begin{lem}[\textbf{Nonlive $\bm{(p_0,x_0)}$ with large $\bm{x_0}$
	reaches a  quasi-dead
	marking $\bm{(p_0,x)}$}]\label{lem:reachquasiPNS}
	There exists a 2-exp function
	$f\qd$ such that 
	for every reversible $d$-dim PNS $G=G^A_{(I,Q)}$ and 
	every nonlive marking $(p_0,x_0)\in Q\times \setN^J$, where
	$J=[1,d]\smallsetminus I$ and
	\begin{equation}\label{eq:condxzero}
x_0\in\upC{\left(2\cdot(|Q|-1)\cdot\nnorm{A}+
		\max\{C\suf^G,\nnorm{A}\cdot C\qd^G\}\right)},
	\end{equation}
		there exists an $f\qd$-quasi-dead marking
	$(p_0,x)\in\Reach((p_0,x_0))$.
\end{lem}
\begin{proof}
	Let $G=G^A_{(I,Q)}$ be a reversible $d$-dim PNS, where
	$J=[1,d]\smallsetminus I$, and let $(p_0,x_0)$
be a nonlive marking with $x_0\in\setN^J$ satisfying	
	\begin{equation}\label{eq:xzerobegin}
		x_0\in\upC{\left(2\cdot(|Q|-1)\cdot\nnorm{A}+C\suf^G\right)}.
	\end{equation}
	The value $C\qd^G=f\qd(\nnorm{G},d)$ will arise from the
	following analysis.

Recall that $Q$ is a
	bSCC in $A\restr{I}$.
If there exists $p\in Q$ that is nonlive in $A\restr{I}$---that is, there exists an action
	$a\in A$ that is not enabled at any $q\in Q$ in
	$A\restr{I}$---then all $q\in Q$ are dead markings in
	$A\restr{I}$ (by our terminology);
consequently, $(p_0,x_0)$ is dead in the net $A$, and thus in the PNS $G$ as well.
	We thus further assume that $Q$ is a live bottom SCC in
	$A\restr{I}$.

	We now fix a dead marking $(q,y)$ such that $(p_0,x_0)\reach
	(q,y)$. 
Let $\emptyset\neq J'\subseteq J$ be such that 
	$y(i)< C\dead^A$ for all $i\in J'$ and 
	\begin{equation}\label{eq:Jprime}
		\textnormal{$(q,y\restr{J'})$ is
		dead in the net $A\restr{(I\cup J')}$.}
	\end{equation}
	By reversibility of $G$ and the
	assumption~\eqref{eq:xzerobegin},
 we may fix an
	execution 
\begin{equation*}
(p_0,x_0)\step{\sigma_0}(q,y')\reach (q,y)
\end{equation*}
where  $p_0\step{\sigma_0}q$ is a shortest path from $p_0$ to $q$ in the
	graph $(Q,E)$ associated with $G$. 
	(Recall Proposition~\ref{prop:quickreachcycle},
	Lemma~\ref{lem:pnsreachfirst}, and the fact that $C\suf^G\geq
	C\vr^G$ by~\eqref{eq:CsufG}.)
Since $|\sigma_0|\leq |Q|-1$, we obtain that, for all $i\in J$,
	\begin{equation}\label{eq:yprimefirst}
		y'(i)\geq x_0(i)-(|Q|-1)\cdot\nnorm{A},
	\end{equation}
and thus
	\begin{equation}\label{eq:yprime}
		y'\in\upC{\left((|Q|-1)\cdot\nnorm{A}+C\suf^G\right)}.
	\end{equation}
Since $y'\vreach y$ in the
	$|J|$-dim ``simple-cycle'' net $\ASC^G$ associated with $G$
	(by Corollary~\ref{cor:cycles}), 
due to the reversibility of $G$---implying the reversibility of
	$\ASC^G$---we may consider a segmented
 virtual execution  
	\begin{equation}\label{eq:downupPNS}
y=y_0\vreach y_1\vreach y_2\cdots\vreach y_k=y'
	\end{equation}
	as in~(\ref{eq:virtsegment}), for the net $\ASC^G$.	
	Consequently, for all $j\in[0,k-1]$ we have
	\begin{equation}\label{eq:CshiftASC}
	\nnorm{y_{j+1}-y_j}\leq
		C\shift^{\ASC^G}=f\shift(\nnorm{\ASC^G},|J|).
\end{equation}
Moreover, 
 the following holds for each $i\in J$:
\begin{itemize}
\item If $y(i)<y'(i)$, then
	$y(i)<y_1(i)<y_2(i)\cdots <y_{j_i}(i)=y_{j_i+1}(i)\cdots
	=y_k(i)=y'(i)$;
\item	
If $y(i)\geq y'(i)$, then $y_j(i)\geq y'(i)$ for all
	$j\in[0,k]$.
	\end{itemize}
By~\eqref{eq:yprime}, 
	there exists $j\leq C\suf^G+(|Q|-1)\cdot\nnorm{A}$ such that 
	the marking $y_j$ in~\eqref{eq:downupPNS}
	is above the level $C\suf^G+(|Q|-1)\cdot\nnorm{A}$. Setting $x'=y_j$, we
	obtain
 \[(q,y)\vreach (q,x')\step{\sigma'}(p_0,x),\]
where
$q\step{\sigma'}p_0$ is a shortest path from $q$ to $p_0$ in the graph
$(Q,E)$. Consequently, $|\sigma'|\leq |Q|-1$, and $x',x$ are both above
$C\suf^G$, because
\begin{equation}\label{eq:xprimex}
	\nnorm{x'-x}\leq (|Q|-1)\cdot\nnorm{A}.
\end{equation}
Since $(p_0,x_0)\reach(q,y)\vreach(q,x')\step{\sigma'}(p_0,x)$
and all $x_0,x',x$ are above $C\suf^G\geq C\vr^G$,
 we obtain
\begin{equation*}
(p_0,x_0)\reach(q,x')\reach(p_0,x)\reach(p_0,x_0),
\end{equation*}
where the last relation follows from reversibility.
Since $(p_0,x_0)\reach(q,y)$, we obtain
\[(p_0,x_0)\reach(p_0,x)\step{\sigma_1}(q,x')\step{\sigma_2}(q,y)\]
for some action sequences
$\sigma_1,\sigma_2$.
The marking $(p_0,x)$ is a candidate for a quasi-dead marking with
the control state $p_0$.
We can bound the length of $\sigma_1$, using
Lemma~\ref{lem:pnsreachfirst} and~\eqref{eq:xprimex}: 
\begin{equation}\label{eq:sigmaone}
	|\sigma_1|\leq C\vr^G\cdot \nnorm{x'-x}\leq
C\vr^G\cdot(|Q|-1)\cdot\nnorm{A}.
\end{equation}
However, we have no bound on $|\sigma_2|$. To handle this, we proceed
as in the proof of Lemma~\ref{lem:reachquasi}.
We set 
\begin{equation}\label{eq:Jprimeprime}
	J''=\{i\in J\mid y(i)< y'(i)\}=\{i\in J\mid y(i)<x'(i)\},
\end{equation}
	noting that $J'\subseteq J''\subseteq J$; recall~\eqref{eq:Jprime} and
	the fact that 
		 $y(i)<C\dead^A\leq C\suf^G\leq y'(i)$ for
	all $i\in J'$.
Since the marking  $(q,y\restr{J'})$ is dead in the net  $A\restr{(I\cup
J')}$, the marking
$(q,y\restr{J''})$ is dead in the net $A\restr{(I\cup J'')}$.

Since $(q,x'\restr{J''})\reach(q,y\restr{J''})$ in  
 $A\restr{(I\cup J'')}$, and  $(q,x'\restr{J''})\geq
 (q,y\restr{J''})$,
 we may apply
 Proposition~\ref{prop:simpletrick}. Consequently, there exists 
 $x''\in \setN^{J''}$ such that
 $x''\geq x'\restr{J''}$ 
and $(q,x'')\step{\sigma}(q,y\restr{J''})$ for some sequence $\sigma$,
where both $\nnorm{x''-x'\restr{J''}}$ and $|\sigma|$ are bounded by
\begin{equation}\label{eq:boundfmr} 
f\mr(\nnorm{A},d)\cdot \nnorm{x'-y}\leq C\mr^A\cdot
	(C\suf^G+(|Q|-1)\cdot\nnorm{A})\cdot C\shift^{\ASC^G}.
\end{equation}
Using reversibility, we obtain $(q,x'\restr{J''})\vreach (q,x'')$
in the net $A'=A\restr{(I\cup J'')}$, as well as in the PNS
$G'=G\restr{(I\cup J'')}=G^{A'}_{(I,Q)}$. Since $x'\restr{J''}$ and
$x''$ are above $C\suf^G\geq C\vr^G\geq C\vr^{G'}$, there exists an
action sequence $\rho$ satisfying
$(q,x'\restr{J''})\step{\rho}(q,x'')$ and 

\begin{equation}\label{eq:boundrho}
	|\rho|\leq C\vr^G\cdot \nnorm{x''-x'\restr{J''}}\leq
C\vr^G\cdot f\mr(\nnorm{A},d)\cdot \nnorm{x'-y}.
\end{equation}
We have thus obtained that, in the net $A\restr{(I\cup J'')}$,
\[(p_0,x\restr{J''})\step{\sigma_1}(q,x'\restr{J''})\step{\rho}(q,x'')\step{\sigma}(q,y\restr{J''}).\]
Inspecting the bound on the length $|\sigma_1\rho\sigma|$ provided
by~\eqref{eq:sigmaone}, \eqref{eq:boundrho}, 
\eqref{eq:boundfmr}, \eqref{eq:CshiftASC},
and \eqref{eq:normASC}, we deduce the existence 
 of a 2-exp function $f\qd$ such that
\[(p_0,x\restr{J''})\step{\sigma_1\rho\sigma}(q,y\restr{J''})\text{ with }
|\sigma_1\rho\sigma|\leq f\qd(\nnorm{G},d).\]
We have $(p_0,x)\step{\sigma_1}(q,x')$ in $G$; however, the sequence $\rho\sigma$ might not be executable from $(q,x')$ in
$G$.
Recalling~\eqref{eq:Jprimeprime}, \eqref{eq:yprimefirst}, and the definition 
 $x'=y_j$ for a respective $y_j$ from~\eqref{eq:downupPNS}, we note
 that,
for all $i\in J\smallsetminus J''$,
\[x'(i)\geq y'(i)\geq x_0(i)-(|Q|-1)\cdot\nnorm{A}.\]
Consequently, if $x_0$ (satisfying~\eqref{eq:xzerobegin}) is above the level
$(|Q|-1)\cdot\nnorm{A}+\nnorm{A}\cdot C\qd^G$, then we obtain 
\[(p_0,x)\step{\sigma_1}(q,x')\step{\rho\sigma}(q,y'')\]
in $G$, where ${y''}\restr{J''}=y\restr{J''}$. Since
$(q,y\restr{J''})$ is dead in $A\restr{(I\cup J'')}$, the marking
$(q,y'')$ is dead in $A$ (and thus in $G$). Consequently, $(p_0,x)$ is an $f\qd$-quasi-dead marking in $G$.
\end{proof}	

\paragraph{Guaranteed existence of live markings $\bm{(p_0,x_0)}$ with large
counter values $\bm{x_0}$.}
We now show that from a live bottom SCC $X$ of a $d$-dim net $A$, we can always extract a
PNS $G=G_{(I,Q)}^A$ where $\nnorm{G}=\max\{\nnorm{Q},\nnorm{A}\}$ is
bounded by a 2-exp function and which possesses a live
marking $(p_0,x_0)$ such that $x_0$ satisfies
condition~\eqref{eq:condxzero}. To this end,
we fix a 2-exp function $f\qd$ provided by Lemma~\ref{lem:reachquasiPNS}.
By $f\suf$, we mean a 2-exp function satisfying
$f\suf(\nnorm{G},d)=C\suf^G=\max\{C\dead^{A},C\vr^G\}$ for every $d$-dim
PNS $G_{(I,Q)}^A$ (which coincides with $A$ if $I=\emptyset$).


\begin{defi}[\textbf{Feasible PNS $\bm{G^{(A,\lambda')}_{(I,Q)}}$ extracted
	from a live bottom SCC $\bm{X}$}]\label{def:feasiblePNS}
	Let a partial function $\bar{\lambda}\colon (\setN_+)^3\to
	\setN_+$ be defined inductively as
	follows for all triples $(m,d,i)\in (\setN_+)^3$ with
	$i\in[1,d]$;
we write $\bar{\lambda}^{(m,d)}_i$ to denote $\bar{\lambda}(m,d,i)$: 
	\begin{itemize}
		\item			
			$\bar{\lambda}^{(m,d)}_1=1+\max\{f\suf(m,d),m\cdot
			f\qd(m,d)\}$;
\item
	$\bar{\lambda}^{(m,d)}_{i+1}=\left[2\cdot((\bar{\lambda}^{(m,d)}_i)^i-1)\cdot
			m+ \max\{f\suf(\bar{\lambda}^{(m,d)}_i,d), m\cdot
			f\qd(\bar{\lambda}^{(m,d)}_i,d)\}\right]+m\cdot
			i\cdot(\bar{\lambda}^{(m,d)}_i)^i.$
	\end{itemize}
For any live bottom SCC $X$ in the reachability graph of
	a $d$-dim net $A$, 
	a $d$-dim PNS $G=G_{(I,Q)}^A$ is \emph{feasible} if 
	$G=G_{(I,Q)}^{(A,\lambda')}$, where
	$\lambda'=(\lambda'_1,\lambda'_2,\dots,\lambda'_d)$ is the extractor satisfying 
	$\lambda'_i=\bar{\lambda}^{(\nnorm{A},d)}_i$ for all
	$i\in[1,d]$. 
	\\
	(Recall Definition~\ref{def:extrPNS} and note that $\lambda'$
	is strictly $\nnorm{A}$-adapted.)
\end{defi}

% Fix Indentation
\noindent
Observe that the function $\bar{f}(m,d)=\bar{\lambda}^{(m,d)}_d$ is a
2-exp function ($\bar{f}(m,d)\leq (2+m)^{2^{\text{poly}(d)}}$).

\begin{prop}[\textbf{Feasible PNSs possess live markings with
	large counter values}]\label{prop:feasiblelive}
	Let $X$ be a live bottom SCC of
	a $d$-dim net $A$, 
and let $G=G_{(I,Q)}^{(A,\lambda')}$ be a feasible PNS extracted from
	$X$.
Then there exists  a (live)	marking
	$(p_0,x_0)\in X$ with $p_0\in Q=X\restr{I}$ and
	\[x_0\in\upC{\left(
2\cdot(|Q|-1)\cdot\nnorm{A}+
		\max\{C\suf^G,\nnorm{A}\cdot C\qd^G\}\right)}.\]
\end{prop}	
\begin{proof}
Let $A,X,G=G_{(I,Q)}^{(A,\lambda')}$ be as in the statement.

		If $I=[1,d]$, then $Q=X$, and the claim holds
		trivially.
	If $I=\emptyset$, then the PNS $G$ coincides with the net $A$
	(in this case, $Q=\{()\}$ and thus $|Q|=1$), and, by
	definition of $G_{(I,Q)}^{(A,\lambda')}$, 
	there exists $x_0\in X$ such that, for all
	$i\in[1,d]$,
	\[x_0(i)\geq 	\lambda'_1\geq
	1+\max\{f\suf(\nnorm{A},d),\nnorm{A}\cdot
			f\qd(\nnorm{A},d)\}\geq
 \max\{C\suf^A,\nnorm{A}\cdot
	C\qd^A\}.\]
	We therefore further assume $\emptyset\neq I\neq [1,d]$.
By definition of $G=G_{(I,Q)}^{(A,\lambda')}$,
we have $\nnorm{Q}<\lambda'_{|I|}$, and there exists $(p_0,x_0)\in X$
	where, for every
	$i\in J=[1,d]\smallsetminus I$, 
\[x_0(i)\geq\lambda'_{|I|+1}-\nnorm{A}\cdot
	|I|\cdot(\lambda'_{|I|})^{|I|}.\]
	The definition of $\lambda'$ (in
	Definition~\ref{def:feasiblePNS}) implies that,  for every
	$i\in J=[1,d]\smallsetminus I$, 
	\[x_0(i)\geq
	2\cdot((\lambda'_{|I|})^{|I|}-1)\cdot\nnorm{A}+
		\max\{f\suf(\lambda'_{|I|},d),\nnorm{A}\cdot
		f\qd(\lambda'_{|I|},d)\}.\]
	Since $(\lambda'_{|I|})^{|I|}-1\geq (1+\nnorm{Q})^{|I|}-1\geq
	|Q|-1$ and $\lambda'_{|I|}\geq
	\max\{\nnorm{Q},\nnorm{A}\}=\nnorm{G}$,
	we obtain,  for every
	$i\in J=[1,d]\smallsetminus I$, 
	\[x_0(i)\geq
	2\cdot(|Q|-1)\cdot\nnorm{A}+
		\max\{f\suf(\nnorm{G},d),\nnorm{A}\cdot
		f\qd(\nnorm{G},d)\},\]
that is,
	\[x_0(i)\geq
2\cdot(|Q|-1)\cdot\nnorm{A}+
		\max\{C\suf^G,\nnorm{A}\cdot C\qd^G\}.\qedhere\]
\end{proof}	


\paragraph{Live solutions of $\bm{(G,(p_0,x_0),\lambda)}$-linear systems.}
Recall Definition~\ref{def:crucialLSPNS} of $(G,(p_0,x_0),\lambda)$-linear systems.
We now show that, given a feasible PNS $G=G^{(A,\lambda')}_{(I,Q)}$,
if $(p_0,x_0)$ is a live marking of $G$ guaranteed by
Proposition~\ref{prop:feasiblelive}
and $S$ is a $(G,(p_0,x_0),\lambda)$-linear system for an extractor
$\lambda$ analogous to that in Lemma~\ref{lem:finalbasicnets},
then $(p_0,\bar{x}_0)$ is a live marking of $G$ 
for every solution
$(\bar{x}_0,(\bar{x}_{J'})_{\emptyset\neq J'\subseteq J})$
of $S$.
The proof is analogous to that of Lemma~\ref{lem:finalbasicnets} and
is therefore presented here in a more concise manner.


\begin{lem}[\textbf{$\bm{(G,(p_0,x_0),\lambda)}$-solutions are live,
	if $\bm{G}$ is feasible and 
	 $\bm{x_0,\lambda}$ are large}]\label{lem:finalPNS}
	Let $G=G_{(I,Q)}^{(A,\lambda')}$ be a feasible PNS, extracted
	from a live bottom SCC $X$ of
	a $d$-dim net $A$.
Let
	$\lambda=(\lambda_1,\lambda_2,\dots,\lambda_{|J|})$,
	where $J=[1,d]\smallsetminus I$, be an extractor
	satisfying:
	\begin{itemize}
		\item
		$\lambda_1\geq
C\suf^G+\nnorm{A}\cdot C\qd^G$;
\item
	$\lambda_{i+1}\geq \lambda_1 + \nnorm{A}\cdot C\vr^{G}\cdot
			(\lambda_i-C\suf^G)$ for each
			$i\in[1,|J|{-}1]$.
	\end{itemize}
	Let $(p_0,x_0)$ be a live marking of $G$ where $x_0\in\setN^J$
	satisfies 
	\[x_0\in\upC{\left(
2\cdot(|Q|-1)\cdot\nnorm{A}+
		\max\{C\suf^G,\nnorm{A}\cdot C\qd^G\}\right)}.\]
Let $S$ be a $(G,(p_0,x_0),\lambda)$-linear system,
and let $(\bar{x}_0,(\bar{x}_{J'})_{\emptyset\neq
	J'\subseteq J})$ be a solution of $S$.
Then
 $(p_0,\bar{x}_0)$ is a live marking. 
\end{lem}	
\begin{proof}
Let $A,G=G_{(I,Q)}^{(A,\lambda')},\lambda$, and $(p_0,x_0)$ be as in
	the statement.
Let $A'=\ASC^G$ be the ``simple-cycle'' $|J|$-dim net associated with
	$G$, where $J=[1,d]\smallsetminus I$.

	Let $S$ be a $(G,(p_0,x_0),\lambda)$-linear
	system~\eqref{eq:GxLS}, namely
\begin{equation*}
	(\bvar{x}\equiv_B x_0)
	\land
	\bigwedge_{\emptyset\neq J'\subseteq
	J}\big((\bvar{x}_{J'}\equiv_B x^\close_{J'})\land
	(\bvar{x}\virtual{*}\bvar{x}_{J'})\big),
\end{equation*}
	where $B=\max\{C\suf^G,\lambda_{|J|}+C\shift^{A'}\}$
and
	$\bvar{x}\virtual{*}\bvar{x}_{J'}$ refers to the virtual
	reachability in $A'$. 
	Recall that $(p_0,x_0)\reach (p_0,x^\close_{J'})$ for each
	$\emptyset\neq J'\subseteq J$; this implies $x_0\vreach
	x^\close_{J'}$ in $A'$.

Towards a contradiction, suppose that
	\begin{center}
		$(\bar{x}_0,(\bar{x}_{J'})_{\emptyset\neq J'\subseteq
		J})$ is a solution of $S$
		where $(p_0,\bar{x}_0)$ is nonlive. 
	\end{center}
	The sets of indices of $\lambda$-small components 
	$I^{\lambda}_{\bar{x}_{J'}}\subseteq J$ and
	$I^{\lambda}_{x^\close_{J'}}\subseteq J$
	(as defined in	Proposition~\ref{prop:Ilambdax}) are
	identical,
	and the vectors $\bar{x}_{J'}$ and $x^\close_{J'}$---viewed as
	markings of the net $A'=\ASC^G$---coincide on
	these indices.
	By definition, all vectors $x_0,x^\close_{J'}$, and therefore also 
	$\bar{x}_0, \bar{x}_{J'}$, are above the level $C\suf^G$.
	Lemma~\ref{lem:reachquasiPNS} implies the existence of
	\begin{center}
	a quasi-dead
		marking $(p_0,x)\in\Reach((p_0,\bar{x}_0))$ with
		$x\in\upC{C\suf^G}$; let
		$I'=I^{\lambda}_{x}\subseteq J$. 
	\end{center}		
	Since $(p_0,x)$ is a quasi-dead marking, we have
	$(p_0,x)\step{\sigma_1}(q,y)$ where
	$(q,y)$ is a dead marking and
	$|\sigma_1|\leq C\qd^G$.  
Note that $I'\neq\emptyset$.
	Indeed, we have $y(i)< C\dead^A$
	for at least one $i\in J$, since otherwise $q$ would be dead
	in $A\restr{I}$, which would imply that $(p_0,x_0)$ is
	nonlive (because
	$(p_0,x_0)\step{\sigma_1}(q,x_0+\Delta(\sigma_1)\restr{J})$). 
	The case $I'=I^{\lambda}_{x}=\emptyset$ would imply
that $x$ is above the level $\lambda_1\geq C\suf^G+\nnorm{A}\cdot
	C\qd^G\geq C\suf^G+\nnorm{A}\cdot |\sigma_1|$, in which case $y$ would be above $C\suf^G\geq
	C\dead^A$---a contradiction. Consequently, 
	$\bar{x}_{I'}$ is well-defined.

	Since $(p_0,\bar{x}_0)\reach (p_0,x)$, we have
	$\bar{x}_0\vreach x$ in $A'$. By reversibility of $A'$
	(both $A$ and $G$ are reversible), it follows that
	$\bar{x}_{I'}\vreach x$ in $A'$ (since $\bar{x}_0\vreach
	\bar{x}_{I'}$ and $\bar{x}_0\vreach x$); this implies
	$(p_0,\bar{x}_{I'})\reach (p_0,x)$ in $G$ (and in $A$), because
	both $\bar{x}_{I'}$ and $x$ are above $C\suf^G$.

The equalities 
\[\dist{\lambda}{I'}{x}=\dist{\lambda}{I'}{\bar{x}_{I'}}=\dist{\lambda}{I'}{x^\close_{I'}}=0\] 
are established analogously to Lemma~\ref{lem:finalbasicnets},
by considering the distance of $\bar{x}_{I'}$ to the $\lambda$-case $I'$
	and a segmented virtual execution $\bar{x}_{I'}\vreach x$ in
	$A'=\ASC^G$.

For each $i\in I'$, the values $x(i)$ and
$\bar{x}_{I'}(i)=x^\close_{I'}(i)$ are in the interval
$[C\suf^G,\lambda_{|I'|}-1]$; therefore, 
$\nnorm{x\restr{I'}-(\bar{x}_{I'})\restr{I'}}\leq \lambda_{|I'|}-C\suf^G$.
The fact 
$(p_0,\bar{x}_{I'})\reach (p_0,x)$ implies that
there exists an action sequence $\sigma_2$ such that 
$(p_0,(\bar{x}_{I'})\restr{I'})\step{\sigma_2} (p_0,x\restr{I'})$ in
the restricted PNS $G\restr{I\cup I'}$ (that is, in $A\restr{I\cup I'}$)
and
\[|\sigma_2|\leq f\vr(\nnorm{G},|I\cup I'|)\cdot
\nnorm{x\restr{I'}-(\bar{x}_{I'})\restr{I'}}\leq C\vr^{G}\cdot
(\lambda_{|I'|}-C\suf^G).\]
Combining with $(p_0,x)\step{\sigma_1}(q,y)$, where $|\sigma_1|\leq
C\qd^G$ and $(q,y)$ is dead, we obtain
\[(p_0,(\bar{x}_{I'})\restr{I'})=(p_0,(x^\close_{I'})\restr{I'})\step{\sigma_2}
(p_0,x\restr{I'})\step{\sigma_1}(q,y\restr{I'})\]
in $G\restr{I\cup I'}$, where $|\sigma_2\sigma_1|\leq C\vr^{G}\cdot
(\lambda_{|I'|}-C\suf^G)+C\qd^G$.
Since for each $i\in J\smallsetminus I'$, we have
$x^\close_{I'}(i)\geq\lambda_{|I'|+1}\geq \nnorm{A}\cdot
|\sigma_2\sigma_1|$, it follows that 
\[(p_0,x^\close_{I'})\step{\sigma_2\sigma_1}(q,y')\]
is an execution in $G$, where $y'\restr{I'}=y\restr{I'}$.
We complete the proof by showing that $(q,y')$ is dead, which
contradicts the assumption that $(p_0,x^\close_{I'})$ is live (since 
$(p_0,x_0)\reach(p_0,x^\close_{I'})$ and $(p_0,x_0)$ is live).

Recall that $(p_0,x)\step{\sigma_1}(q,y)$ and $|\sigma_1|\leq C\qd^G$.
Since for each $i\in J\smallsetminus I'$, we have $x(i)\geq
\lambda_{|I'|+1}\geq \lambda_1\geq C\dead^A+ \nnorm{A}\cdot C\qd^G$, it follows that 
$y(i)\geq C\dead^A$ for each  $i\in J\smallsetminus I'$. This implies
that $(q,y\restr{I'})$ is dead in $A\restr{I\cup I'}$, which in turn
implies that $(q,y')$ is dead in $A$ (since
$y'\restr{I'}=y\restr{I'}$).
\end{proof}

\paragraph{Proof of Theorem~\ref{th:upper}.} 

Given a $d$-dim
net $A$ possessing a live bottom SCC $X\subseteq\setN^d$, we consider 
 a feasible PNS  $G=G_{(I,Q)}^{(A,\lambda')}$ as defined in
Definition~\ref{def:feasiblePNS}. We set $J=[1,d]\smallsetminus I$.

Since $\nnorm{Q}<\lambda'_{|I|}$, it
follows that 
\[\nnorm{G}=\max\{\nnorm{Q},\nnorm{A}\}\leq
(2+\nnorm{A})^{2^{\text{poly}(d)}}.\]
The 2-exp bound follows from the definition of $\lambda'$ and
the closure properties of 1-exp and 2-exp functions discussed at the
beginning of this section.
Moreover, by Proposition~\ref{prop:feasiblelive}, there exists a live	marking
	$(p_0,x_0)\in X$ with $p_0\in Q=X\restr{I}$ and
	\[x_0\in\upC{\left(
2\cdot(|Q|-1)\cdot\nnorm{A}+
		\max\{C\suf^G,\nnorm{A}\cdot C\qd^G\}\right)}.\]
Recall that $|Q|\leq (1+\nnorm{Q})^{|I|}$ (where $|I|\leq d$) and
that $f\suf$ and $f\qd$ are 2-exp functions.
Lemma~\ref{lem:finalPNS} implies that every solution
of a $(G,(p_0,x_0),\lambda)$-linear system $S$, 
where $\lambda$ is defined
by 
\[\lambda_1=
C\suf^G+\nnorm{A}\cdot C\qd^G, \text{ and } 
\lambda_{i+1}= \lambda_1 + \nnorm{A}\cdot C\vr^{G}\cdot
			(\lambda_i-C\suf^G)\text{ for each }
			i\in[1,|J|{-}1],\]
			yields  
a live marking $(p_0,\bar{x}_0)$.
The dimension of $S$ is at most $d\cdot 2^d$, and
\[\max\{\mlcm{S},\nnorm{S}\}\leq (2+\nnorm{A})^{2^{\text{poly}(d)}}.\]
Indeed, the virtual reachability in $S$ refers to the ``simple-cycle
net''  $\ASC^G$; consequently,
Corollary~\ref{cor:virtual2system} implies that 
\[\mlcm{S}\leq |J|\,!\cdot\nnorm{\ASC^G}\leq d\,!\cdot
|Q|\cdot\nnorm{A}\leq  d\,!\cdot (1+\nnorm{Q})^d \cdot\nnorm{A},\]
where $\nnorm{Q}\leq (2+\nnorm{A})^{2^{\text{poly}(d)}}$.
Since
$B=B_{(G,\lambda)}=\max\{C\suf^G,\lambda_{|J|}+C\shift^{\ASC^G}\}$,
where
\begin{equation*}
\lambda_{|J|}\leq (2+\nnorm{A})^{2^{\text{poly}(|J|)}}\leq (2+\nnorm{A})^{2^{\text{poly}(d)}}
	\text{ and }f\shift(\nnorm{\ASC^G},|J|)\leq f\shift(|Q|\cdot \nnorm{A},d),
\end{equation*}
	we also obtain $\nnorm{S}\leq  (2+\nnorm{A})^{2^{\text{poly}(d)}}$.

Analogously to Proposition~\ref{prop:Affsmall}, we use
Theorem~\ref{thm:smallsolutions} to deduce that
there exists a solution of $S$ that yields $(p_0,\bar{x}_0)$ (which is
a live marking of $A$) such that
\[\nnorm{(p_0,\bar{x}_0)}\leq
(2+\nnorm{A})^{2^{\text{poly}(d)}}.\tag*{\qed}\]


\subsection{Dead markings}\label{subsec:dead}
\hfill\\

In this section, we provide further details regarding the
characterisation~\eqref{eq:alldead} of dead markings.
Recall that in our terminology, a marking $y\in\setN^d$ of a $d$-dim net $A$
is dead if there exists an action $a\in A$ that is dead at $y$, that is,
$a$ is disabled at every marking $y'\in\Reach(y)$. The following
proposition is a straightforward observation.
\begin{prop}[\textbf{If a restriction of a marking is dead, then the
	marking is dead}]\label{prop:restrdead}
	Given a $d$-dim net $A$ and $I\subseteq [1,d]$, if 
 $y\restr{I}$ is dead in $A\restr{I}$, then $y$ is dead in $A$.
\end{prop}
\begin{proof}
The claim follows from the fact that, for every $y\in\setN^d$, the set
$\{y'\restr{I}\mid y\reach y'$ in $A\}$ is a subset of the set
$\{y'\in\setN^I\mid y\restr{I}\reach y'$ in $A\restr{I}\}$.
\end{proof}

% Fix Indentation
\noindent
We now adapt Rackoff's result for coverability
from~\cite{DBLP:journals/tcs/Rackoff78}.
\begin{lem}[\textbf{Deadness is determined by components with
	``small'' values}]\label{lem:rackoffdead}
\hfill\\
	There exists a 2-exp function $f\dead$ with the following property:
\\
For every
$d$-dim net $A$, a marking $y\in\setN^d$ is dead in $A$ if and only if 
 $y\restr{\textsc{s}_y}$ is dead in the restricted net
	$A\restr{\textsc{s}_y}$,
	where $\textsc{s}_y=\{i\in [1,d]\mid
y(i)<f\dead(\nnorm{A},d)\}$.
\end{lem}
\begin{proof}
Note that an action $a=(a_-,a_+)$ in a net $A$
	is not dead
	at a marking $x$ if and only if $x$ can cover $a_-$, that is, 
	if and only if there exists $x'$ such that $x\reach x'\geq a_-$.
	Along the lines of~\cite{DBLP:journals/tcs/Rackoff78}
	(Lemmas~3.3. and~3.4), we define a function
	$g:\setN\times\setN\to\setN$ as follows:
\begin{itemize}
	\item			
 $g(m,0)=0$;
\item
	for $i\geq 1$,	$g(m,i)= (m + m\cdot g(m,i{-}1))^{i}+g(m,i{-}1)$.
\end{itemize}		

	% Fix Indentation
	\noindent
	By induction on $|I|$, it is straightforward to verify that 
	for a $d$-dim net $A$, 
an action $a=(a_-,a_+)\in A$, an index set $I\subseteq [1,d]$, and a
	marking
	$x\in\setN^I$ of $A\restr{I}$, the following holds: if  $x$ can cover
	$(a_-)\restr{I}$ (in the net $A\restr{I}$), then 
there exist an action sequence $\sigma$ and a marking $x'\in\setN^I$ such that 
	\begin{equation*}
		x\step{\sigma}x'\geq (a_-)\restr{I}\text{ and }|\sigma|\leq
	g(\nnorm{A},|I|). 
	\end{equation*}
The claim is trivial for $I=\emptyset$.	In the case $I\neq\emptyset$,
we observe that 
if a shortest execution  $x\reach x'\geq (a_-)\restr{I}$ in
	$A\restr{I}$
	 has a prefix
	reaching the value $n=\nnorm{A} + \nnorm{A}\cdot
	g(\nnorm{A},|I|{-}1)$ or larger in some
	component, then the length of the suffix is bounded
	by $g(\nnorm{A},|I|{-}1)$ by the inductive hypothesis, and the
	length of the prefix is bounded by $n^{|I|}$.

This implies that for any $d$-dim net $A$ and $x\in\setN^d$, 
the values
	$x(i)$, $i\in[1,d]$, for which  
	$x(i)\geq \nnorm{A}+\nnorm{A}\cdot g(\nnorm{A},d)$ are
	irrelevant to 
	the question of whether $x$ is dead. 
As the required 2-exp function, we can thus take
	\[f\dead(m,i)=m+m\cdot g(m,i).\qedhere\] 
\end{proof}	

\begin{rem}
We note that the recent
paper~\cite{kunnemann_et_al:LIPIcs.ICALP.2023.131} has improved Rackoff's
bound in~\cite{DBLP:journals/tcs/Rackoff78}, but such fine-tuned
results are not needed at our level of analysis.
\end{rem}


\subsection{Virtual reachability in reversible 
nets}\label{subsec:virtrevers}
\hfill\\

In this section, we formalize the previous informal claim regarding small-segment
virtual executions~(\ref{eq:virtsegment})
that are ``directed'' from a marking $x$ to a marking $y$.


We recall the notation $A\D=\{\Delta(a)\mid a\in
A\}$ and the fact that
the virtual reachability relation $\vreach$ in any
\emph{reversible} $d$-dim net $A$ is
symmetric: $x\virtual{*}y$ implies $y\virtual{*}x$.
This follows from the fact that $x\virtual{*}y$ in $A$ is equivalent to
$(y-x)\in A\D^*$, and that 
$A\D^*$  is a subgroup of
$(\setZ^d,+)$ if $A$ is reversible.


For any $d\in\Nat$, we define the function 
\begin{equation*}
\sign:\setZ^d\to\{-1,0,1\}^d
\end{equation*}
as follows:
 for $x\in\setZ^d$ and $i\in[1,d]$, we set
\begin{itemize}
	\item		
 $\sign(x)(i)=-1$ if $x(i)<0$;
\item
$\sign(x)(i)=0$ if $x(i)=0$;
\item
		$\sign(x)(i)=1$ if $x(i)>0$.
\end{itemize}
		We define the following partial order $\sleq$ on $\setZ^d$:
\begin{equation*}
	\textnormal{$x\sleq y$\ if $\sign(x)=\sign(y)$ and $|x(i)|\leq
	|y(i)|$ for all $i\in [1,d]$.} 
\end{equation*}
For instance, $(0,5,-3)\sleq (0,11,-6)$, but 
 $(0,5,-3)\not\sleq (1,11,-6)$, since
 $\sign((0,5,-3))=(0,1,-1)\neq\sign((1,11,-6))=(1,1,-1)$.
Note that $\mathbf{0}\sleq x$ implies $x=\mathbf{0}$.
For $X\subseteq \setZ^d$, 
\[\minsleq{X}\]
denotes the set of 
minimal elements of $X$ w.r.t.\ $\sleq$\,; it is finite
since
$\sleq$ is a well-partial-order.

\begin{prop}[\textbf{``Directed'' decomposition of elements in group $\bm{L}$ by
	$\bm{\minsleq{L}}$}]\label{prop:onestepsubgroup}
	If $L\subseteq \setZ^d$ is a group and
	$y\in L$, then there exists $z\in\minsleq{L}$ such that
	$z\sleq y$, $y=z+(y-z)$,   
	and $(y-z)\in L$. Moreover, $\nnorm{y-z}<\nnorm{y}$
 if $y\neq\mathbf{0}$.
\end{prop}	
\begin{proof}
Let $L\subseteq \setZ^d$ be a group, and
let $y\in L$.
	Since $\sleq$ is a well-partial-order, there exists $z\in\minsleq{L}$ such that $z\sleq
	y$. We have $(y-z)\in L$ because $L$ is a group.
To show the strict decrease in norm, consider a component $i\in[1,d]$:
	\begin{itemize}
		\item
			If $y(i)=0$, then $z(i)=0$, since
			$\sign(y)=\sign(z)$; therefore, $y(i)-z(i)=0$.
\item
If $y(i)>0$, then $0<z(i)\leq y(i)$, 
which implies $0\leq y(i)-z(i)<y(i)$; therefore,
$|y(i)-z(i)|<|y(i)|$.
\item
If $y(i)<0$, then $y(i)\leq z(i)<0$, 
			which implies $y(i)<y(i)-z(i)\leq 0$;
			therefore, 
$|y(i)-z(i)|<|y(i)|$.
\end{itemize}
Consequently, $\nnorm{y-z}<\nnorm{y}$ if
	$y\neq\mathbf{0}$.
\end{proof}

\begin{cor}[\textbf{Segmenting virtual executions of reversible nets}]\label{cor:subgroup}
\hfill\\
	Given a reversible net $A$ (for which $A\D^*$ is a group), if
	$x\virtual{*}y$, then 
	there exist $z_1,z_2,\dots,z_k$ in $\minsleq{A\D^*}$
such that
\begin{equation}\label{eq:decompseq}
x\virtual{*}(x+z_1)\virtual{*}(x+z_1+z_2)\cdots\virtual{*}(x+z_1+z_2\cdots
	+z_k)=y
\end{equation}	
and for each $j\in[1,k]$ it holds that
 $z_j\sleq y-(x+z_1+z_2\cdots +z_{j-1})$.
\end{cor}
\begin{proof}
This is a direct consequence of Proposition~\ref{prop:onestepsubgroup},
since $x\virtual{*}y$ is equivalent to $(y-x)\in A\D^*$. 
By applying the proposition repeatedly, we construct the sequence:
in the first step, we choose
$z_1\in\minsleq{A\D^*}$ such that 
$z_1\sleq (y-x)$. Thus $y-x=z_1+(y-x-z_1)$, which implies
$y-x-z_1=y-(x+z_1)\in A\D^*$; that is,
$(x+z_1)\vreach y$. Moreover, if  $x\neq y$, then
$\nnorm{y-(x+z_1)}<\nnorm{y-x}$.
\end{proof}

% Fix Indentation
\noindent
Recall that we used
$C\shift^A=f\shift(\nnorm{A},d)$ to bound the norms of segments in the virtual
executions~(\ref{eq:virtsegment}) of a $d$-dim net $A$. In principle,
by~\eqref{eq:Cshift}
 we claimed  
that $\nnorm{\minsleq{A\D^*}}\leq f\shift(\nnorm{A},d)$, where
$f\shift$ is a 1-exp function. The following lemma establishes this
result; it is slightly more general, since it shows that 
	\[\nnorm{\minsleq{A^*}}\leq (2+\nnorm{A})^{\text{poly}(d)}\] 
for every finite set $A\subseteq \setZ^d$ (where $A^*$ is not necessarily a group).

\begin{lem}[\textbf{For every finite set $\bm{A\subseteq \setZ^d}$,	
the norm $\bm{\nnorm{\minsleq{A^*}}}$ is 1-exp}]\label{lem:Lminbound}
\hfill\\
For every finite set $A\subseteq \setZ^d$, it holds that
	\[\norm{}{\minsleq{A^*}}\leq  (3+\nnorm{A}\cdot |A|)^d\leq
	(3+\nnorm{A}\cdot (1+2\cdot\nnorm{A})^d)^d.\]
\end{lem}
\begin{proof}
By decomposing the set $A^*$ according to the function $\sign$
 (recall that $x\sleq y$ requires $\sign(x)=\sign(y)$, so vectors in
 different sign classes are incomparable), we obtain:
		\[\minsleq{A^*}=\bigcup_{s\in \{-1,0,1\}^d}\minsleq{\{y\in A^*
	\mid \sign(y)=s\}}.\]
	Thus, it suffices to provide the required bound on
	$\nnorm{\minsleq{\{y\in A^* \mid \sign(y)=s\}}}$ 
	for an arbitrary vector  $s\in\{-1,0,1\}^d$, which is henceforth fixed.
Observe that by flipping the sign of the $i$th component of each
	vector in $A$ for all $i\in[1,d]$ such that $s(i)=-1$, we may
	assume without loss of generality that $s\in \{0,1\}^d$. We
	set  \[I=\{i\in[1,d]\mid s(i)=0\}.\]
	For $k=|A|$, we henceforth view $A$ as a $d\times k$ matrix
	(with $d$ rows and $k$ columns), where
the columns correspond to the elements of $A$.
	Consequently, we now aim to provide a bound on the
	norm $\nnorm{\min{}_{\leq}(Y)}$, where $\leq$
	denotes the
	standard component-wise order and the set $Y\subseteq \setN^d$
	is defined as follows:
	\[Y=\{Av \mid v\in\setN^k
	\land \bigwedge_{i\in I}(Av)(i)=0\land
	\bigwedge_{i\in[1,d]\smallsetminus I}(Av)(i)\geq 1\}.\]
We proceed similarly to the proof
	of~\cite[Corollary~1]{DBLP:conf/rta/Pottier91}, which deals with
systems of linear inequalities.
Let $B$ be the $d \times (k+d+1)$ matrix
\[B=\left(\begin{array}{c|c|c} A & -\mathrm{Id}_d & -s\end{array}\right),\]
that is, the $k$ columns of $A$, followed by the columns of the negated
$d\times d$ identity matrix, followed by the single column $-s$.
We set 
	\[M=\{x\in\setN^{k+d+1}\mid
	Bx=\mathbf{0}\}=\{(v,z,t)\in\setN^k\times\setN^d\times\setN\mid
	Av-z-ts=\mathbf{0}\},\]
where the three blocks of $B$ act on $v$, $z$, and $t$, respectively.
	We then observe that 
	\begin{equation}\label{eq:YviaM}
		Y=\{Av\mid \exists z\in\setN^d: (v,z,1)\in M\land
	\bigwedge_{i\in I}z(i)=0\}.
	\end{equation}
 Lemma~\ref{lem:smallsolution} implies that $M=X^*$, 
		where $X=\min{}_{\leq}(M\smallsetminus\{\mathbf{0}\})$ and,
		since $\rank{B}\leq d$ and $\norm{1}{B}\leq
		k\cdot\norm{}{A}+2$,
	\begin{equation}\label{eq:XboundforM}
		\norm{1}{X}\leq (3+k\cdot\norm{}{A})^d.
	\end{equation}
Finally, we fix an arbitrary  $y\in \min{}_{\leq}(Y)$ and bound its
	norm $\nnorm{y}$.
	By~\eqref{eq:YviaM}, we may fix $v$ and $z$ such that $y=Av$,
	$(v,z,1)\in M$,
and $z(i)=0$ for all $i\in I$.
	Since $M=X^*$, the tuple $(v,z,1)$ is a sum of elements of $X$;
	their last components are nonnegative and sum to $1$, so exactly
	one summand $(v',z',1)$ has last component $1$, and it satisfies
	$(v',z')\leq (v,z)$ and
	$\norm{1}{(v',z',1)}\leq (3+k\cdot\norm{}{A})^d$
	by~\eqref{eq:XboundforM}.
	As $z'\leq z$ and $z(i)=0$ for all $i\in I$, we have $z'(i)=0$ for
	all $i\in I$, and hence $Av'=s+z'$ belongs to $Y$
	by~\eqref{eq:YviaM}.
	Because 
	\[Av'=s+z'\leq s+z=Av=y\] and $y\in \min{}_{\leq}(Y)$, it
	follows that $z'=z$. 
	Consequently, discarding the last component of $(v',z',1)$,
	\[\nnorm{z}\leq\norm{1}{z}=\norm{1}{z'}\leq (3+k\cdot\norm{}{A})^d-1.\]
	Since $Av-z=y-z=s\in\{0,1\}^d$,
	we conclude that $\nnorm{y}\leq\nnorm{z}+1\leq
	(3+k\cdot\norm{}{A})^d$.
\end{proof}

